% Generated by R/51_manuscript_latex.R from the section drafts in docs/manuscript/.
% Edit the section drafts, then re-run R/50_assemble_manuscript.R and this script.
\documentclass[11pt,a4paper]{article}
\usepackage[utf8]{inputenc}
\usepackage[T1]{fontenc}
\usepackage{lmodern}
\usepackage[margin=2.5cm]{geometry}
\usepackage{microtype}
\usepackage{booktabs}
\usepackage{longtable}
\usepackage{array}
\usepackage{calc}
\usepackage{pdflscape}
\usepackage{caption}
\usepackage{graphicx}
\usepackage{float}
\usepackage{amsmath}
\usepackage{textcomp}
\usepackage{lineno}
\usepackage[hidelinks]{hyperref}
\captionsetup{font=small,labelfont=bf,justification=raggedright,singlelinecheck=false}
\providecommand{\tightlist}{\setlength{\itemsep}{0pt}\setlength{\parskip}{0pt}}
\setlength{\parskip}{0.5em}
\setlength{\parindent}{0pt}
\linespread{1.25}
\title{Missing-data methods and equity estimates of child undernutrition: a methodological study using the 2021 Madagascar Demographic and Health Survey}
\author{[CONFIRM author list and affiliations]}
\date{Draft compiled \today}
\begin{document}
\linenumbers
\maketitle

\section{Abstract}\label{abstract}

\textbf{Background.} Household surveys such as the Demographic and
Health Surveys (DHS) are the main source of subnational estimates of
child undernutrition, and are used to decide which regions and which
socioeconomic groups are targeted first. Anthropometric measurements are
incomplete in every survey, yet the analytical treatment of those
missing values is rarely examined, and its effect on the resulting
equity ordering is essentially unknown.

\textbf{Methods.} We analysed the 2021 Madagascar DHS (6,682 de facto
children aged 0--59 months), comparing four approaches to missing
anthropometry under the survey design: complete-case analysis, inverse
probability weighting (IPW), multiple imputation by predictive mean
matching (MI-PMM) and random-forest multiple imputation (MI-RF). We
estimated stunting and wasting prevalence nationally and for 23 regions
and 5 wealth quintiles, measured cross-method agreement in subgroup
rankings against a benchmark of sampling variability, and ran a
simulation (12 scenarios, 500 replicates each) combining a survey
bootstrap with imposed MCAR, MAR, strong-MAR and MNAR missingness at
10--50\%.

\textbf{Results.} A valid height-for-age z-score was unavailable for
4.2\% of children and a weight-for-height z-score for 3.1\%. Missingness
was strongly patterned (44.6\% where the mother's interview was
incomplete versus 2.4\% where she was interviewed), yet national
estimates differed across methods by at most 0.12 percentage points
(complete-case stunting 39.8\%, 95\% CI 38.0--41.7; wasting 7.7\%,
6.9--8.6). Regional rankings agreed closely across methods (Spearman $\rho$ $\ge$
0.987), far more closely than rankings redrawn from sampling uncertainty
alone (median $\rho$ 0.88--0.93). In simulation, all methods were unbiased
under MCAR and under missingness calibrated to the survey itself; under
strong MAR at 50\% missingness complete-case analysis was biased by $-$3.7
percentage points with 11\% CI coverage, while IPW remained unbiased
with 95\% coverage; under MNAR no method recovered the truth.

\textbf{Conclusions.} Here the choice of missing-data method was
immaterial, and sampling uncertainty, not missing data, limited equity
targeting. That equivalence reflects low, weakly outcome-related
missingness rather than a general result: when missingness tracked
undernutrition risk, only response weighting remained valid.

\textbf{Keywords.} missing data; multiple imputation; inverse
probability weighting; complex surveys; child stunting; health equity;
Demographic and Health Surveys; Madagascar

\section{Background}\label{background}

Child undernutrition remains one of the largest contributors to child
mortality and to lifelong deficits in health, schooling and earnings
\cite{ref1,ref2}. Stunting, or low height-for-age, affects about one in
five children under 5 worldwide, and progress is uneven both between and
within countries \cite{ref3,ref4}. Reducing it is a global commitment
under the World Health Assembly targets and the Sustainable Development
Goals \cite{ref5}. In Madagascar, 40\% of children under 5 are stunted,
among the highest prevalences in the world \textbf{{[}CONFIRM ranking
against the 2023 joint malnutrition estimates{]}}, and regional
prevalence ranges from about 22\% to 52\% \cite{ref6}.

Because national averages hide this heterogeneity, nutrition programmes
are increasingly targeted subnationally, and the evidence for those
decisions comes almost entirely from household surveys, above all the
DHS \cite{ref7}. These surveys measure the height and weight of children
in sampled households and publish prevalence by region and wealth
quintile; the resulting tables are used to rank areas and prioritise
resources. Two features of such estimates deserve more scrutiny than
they usually receive. First, they carry sampling uncertainty that is
substantial in small domains, so the ordering of regions is itself
uncertain. Second, anthropometric data are never complete: children are
absent, carers refuse, and some measurements are implausible and
flagged. Reviews of DHS anthropometry show that incompleteness and
flagging vary appreciably between surveys \cite{ref8,ref9}, and cleaning
rules alone can shift estimated malnutrition prevalence \cite{ref10}.

How missing anthropometry is handled is therefore a methodological
choice with potential consequences for who is identified as worst
affected. In practice the choice is usually implicit: survey reports
present complete-case estimates restricted to children with valid
measurements, which are unbiased only if missingness is unrelated to
nutritional status given the design and the covariates used
\cite{ref11}. Alternatives are well established in the statistical
literature --- inverse probability weighting for nonresponse
\cite{ref12} and multiple imputation, including machine-learning
variants \cite{ref13,ref14,ref15} --- and guidance exists on reporting
them \cite{ref16}. What is missing is evidence on what difference they
make in this setting: for prevalence under a complex survey design, for
the subnational ranking that drives targeting, and under missingness
mechanisms realistic for anthropometry.

Related work has approached parts of this problem separately.
Methodological studies of survey nonparticipation have concentrated on
HIV prevalence, where selection models have been used to correct
estimates for people who decline testing \cite{ref17}, and the
statistical literature on ranking has long warned that league tables of
institutions or areas are far less reliable than their point estimates
suggest \cite{ref18,ref19}. Studies of DHS anthropometry, in turn, have
focused on data quality: how often measurements are missing, flagged or
heaped, and how cleaning rules change prevalence \cite{ref10,ref8,ref9}.
What has not been examined, to our knowledge, is the combination that
matters for programme targeting: whether the analytical treatment of
missing anthropometry changes which subnational groups a survey
identifies as worst affected, and how any such change compares with the
ranking uncertainty already present from sampling.

This study addresses that gap using the 2021 Madagascar DHS. We ask
three questions. First, how much do survey-weighted stunting and wasting
prevalence differ between complete-case analysis, IPW and two forms of
multiple imputation, nationally and by region and wealth quintile?
Second, does the choice of method change the equity ordering --- which
regions and quintiles are identified as worst affected --- and how does
any such change compare with the uncertainty that sampling error alone
places on those rankings? Third, in a simulation in which the truth is
known and the mechanism is controlled, what are the bias, precision and
confidence-interval coverage of each method under MCAR, MAR and MNAR
missingness at rates from 10\% to 50\%? Throughout, we keep estimation
inside the survey design, distinguish structurally inapplicable values
from genuinely missing ones, and validate the pipeline by reproducing
the published survey estimates before any method is compared.

\section{Methods}\label{methods}

\subsection{Data source and survey
design}\label{data-source-and-survey-design}

We used data from the 2021 Madagascar Demographic and Health Survey
(Enquête Démographique et de Santé à Madagascar, EDSMD-V), a nationally
representative household survey conducted by the Institut National de la
Statistique (INSTAT) with technical assistance from ICF through The DHS
Program \cite{ref6}. The survey was designed to produce estimates for
the nation, for urban and rural areas, and for 23 study domains: the 22
administrative regions that existed at survey design (with Vatovavy and
Fitovinany forming a single region), with the capital, Antananarivo, and
the remainder of Analamanga region treated as separate domains
\cite{ref6}.

The sample was a stratified two-stage cluster sample drawn from the 2018
General Population and Housing Census (RGPH-3) frame. Each domain was
divided into urban and rural strata (the capital is urban only), giving
45 sampling strata. In the first stage, 657 enumeration areas were
selected with probability proportional to size; in the second stage, 34
households were selected systematically with equal probability in each
enumeration area, and 20,510 households were interviewed \cite{ref6}. In
a random half of households, all children under 5 years were weighed and
measured \cite{ref6}. Fieldwork, launched in March 2020 and suspended
because of the COVID-19 pandemic, restarted on 3 March 2021 and lasted
143 days \cite{ref6}. \textbf{{[}CONFIRM whether data from the 56
clusters completed in March 2020 are included in the released file and
whether this needs mention.{]}}

\subsection{Study population and
outcomes}\label{study-population-and-outcomes}

The analytic population comprised de facto children aged 0--59 months
(children who slept in the household the night before the interview)
listed in the household member (PR) recode file, the population used for
the survey's published nutrition indicators \cite{ref6,ref20} (n = 6,682
children in 649 primary sampling units and 45 strata).

Height-for-age (HAZ), weight-for-height (WHZ) and weight-for-age
z-scores had been calculated by DHS using the WHO Child Growth Standards
\cite{ref21}. The recode stores z-scores multiplied by 100 and uses
separate codes for values that are missing (not measured) and for
measurements that fall outside plausible limits: 9996 (height out of
plausible limits), 9997 (age out of plausible limits) and 9998
(WHO-flagged z-score). We kept these two types of missingness distinct
throughout. Stunting and severe stunting were defined as HAZ \textless{}
$-$2 and \textless{} $-$3, and wasting and severe wasting as WHZ \textless{}
$-$2 and \textless{} $-$3 \cite{ref21}.

\subsection{Covariates and structural versus genuine
missingness}\label{covariates-and-structural-versus-genuine-missingness}

Auxiliary variables were study domain (hereafter ``region''), place of
residence, household wealth quintile, child's sex, child's age in months
(and in the eight age bands used in the Final Report), household size,
mother's interview status, mother's education (none, primary, secondary
or higher) and mother's age. Mother's age was taken from the mother's
own household roster line.

Two sources of incomplete maternal information were distinguished. For
children whose mother did not live in the household (including children
whose mother had died; n = 569), maternal education and age are not
applicable. These were coded as a separate category (``not in
household''), were never imputed, and were represented in every model
through mother's interview status. Genuine missingness comprised
mother's age for children whose mother lived in the household but could
not be linked to a roster line (mother not de facto, n = 30; incomplete
women's interview, n = 193) and mother's education recorded as ``don't
know'' (n = 6).

\subsection{Complex survey design}\label{complex-survey-design}

All estimates accounted for sampling weights (the DHS household weight
divided by 10$^6$), stratification and clustering. Variances were estimated
by Taylor linearisation \cite{ref22}, with primary sampling units nested
within strata; single-PSU strata were handled by centring at the grand
mean. Subgroup prevalences were estimated as domain estimates, retaining
all sampling units in the variance calculation. Because many subgroup
prevalences are low, 95\% confidence intervals (CIs) were constructed on
the logit scale (delta-method variance) and back-transformed, with a t
reference distribution on the design degrees of freedom (number of PSUs
minus number of strata, 604). Associations between outcome missingness
and covariates were assessed with design-based Rao--Scott tests
\cite{ref23}.

\textbf{Pipeline validation.} Before applying any missing-data method,
we checked that the complete-case, design-based estimates reproduced the
published survey figures. Estimates reproduced all 246 cells of Final
Report Table 11.1 (moderate and severe stunting, wasting and underweight
by age, sex, residence, 23 regions and wealth quintile) to the published
precision, the national unweighted and weighted denominators, and the 53
stunting, wasting and underweight entries of the sampling-error appendix
(national standard errors and $\pm$2 SE intervals for residence and the 23
regions) \cite{ref6}. The analysis pipeline stops automatically if any
of these checks fails.

\subsection{Estimand and missing-data methods
(RQ1)}\label{estimand-and-missing-data-methods-rq1}

The target estimand was the survey-weighted prevalence of each outcome
among all de facto children aged 0--59 months, nationally and by region
and wealth quintile. An outcome was considered missing when no valid
z-score was available, whether because the child was not measured or
because the value was flagged. We compared five approaches, chosen to
span the estimators most often used for incomplete survey outcomes
\cite{ref11}.

\begin{enumerate}
\def\labelenumi{\arabic{enumi}.}
\tightlist
\item
  \textbf{Complete-case analysis (CCA).} Children with a valid z-score,
  analysed with the design weights. This is the definition used in DHS
  reports \cite{ref20}.
\item
  \textbf{CCA with complete covariates.} As above, additionally
  requiring observed mother's age and education, representing the common
  practice of restricting to records complete for all analysis
  variables.
\item
  \textbf{Inverse probability weighting (IPW).} For each index
  separately, the probability of a valid z-score was modelled by
  survey-weighted logistic regression on region, residence, wealth
  quintile, sex, age band, mother's interview status, mother's education
  and household size. Each responding child's design weight was divided
  by the fitted response probability \cite{ref12}. To avoid collinearity
  with mother's interview status, children whose mother was not in the
  household were assigned a constant education level in this model, and
  the six ``don't know'' values were set to the modal category. Variance
  estimation treated the estimated weights as fixed.
\item
  \textbf{Multiple imputation by chained equations with predictive mean
  matching (MI-PMM)} \cite{ref13,ref24}. HAZ, WHZ and mother's age were
  imputed by predictive mean matching \cite{ref25}, and mother's
  education by polytomous logistic regression. Imputation models
  included all other variables, including the survey strata (region and
  residence) and the log design weight, so that the imputation model
  reflected the sampling design \cite{ref26}. Weight-for-age was
  excluded from the imputation model: it is close to a deterministic
  function of HAZ and WHZ and is usually missing jointly with them, and
  its inclusion produced non-mixing chains with drifting variance in
  diagnostic runs. Imputation models for the maternal variables were
  fitted only to children whose mother lived in the household, so that
  structurally inapplicable values were neither used as donors nor
  imputed. We generated 20 imputed datasets with 20 iterations each;
  convergence was assessed from trace plots of chain means and standard
  deviations. Binary outcomes were derived from the imputed z-scores
  after imputation.
\item
  \textbf{Random-forest multiple imputation (MI-RF)} \cite{ref14,ref15}.
  As for MI-PMM, but with every incompletely observed variable imputed
  by random forests (100 trees; ranger implementation \cite{ref27}).
  Unlike single random-forest imputation, this procedure propagates
  imputation uncertainty through multiple imputed datasets.
\end{enumerate}

For both MI approaches, the survey design was applied to each completed
dataset, and estimates were combined with Rubin's rules \cite{ref28} on
the logit scale \cite{ref29}. Inference used Barnard--Rubin small-sample
degrees of freedom \cite{ref30} with the design degrees of freedom as
the complete-data value.

\textbf{Sensitivity analysis.} Because implausible measurements may
reflect measurement error, not genuinely extreme values, we repeated all
analyses with flagged children excluded from the target population,
following the DHS and WHO convention. Exclusions were made separately
for each index: flagged-HAZ children for stunting outcomes and
flagged-WHZ children for wasting outcomes. By construction, CCA
estimates are identical in the main and sensitivity analyses.

\subsection{Equity ranking analysis
(RQ2)}\label{equity-ranking-analysis-rq2}

Within each method, regions and wealth quintiles were ranked by
prevalence of stunting and of wasting (rank 1 = highest). Agreement
between methods was summarised for every pair of methods by Spearman's
$\rho$, Kendall's $\tau$-b, the largest absolute rank shift and, as a
targeting-relevant measure, the overlap of the ``worst-affected'' set:
the five highest-prevalence regions (about the top fifth of 23) and the
highest-prevalence wealth quintile. A subgroup was classified as having
method-dependent worst-affected status if it belonged to this set under
some but not all methods.

Because rankings of survey estimates are themselves uncertain, we
quantified sampling uncertainty in ranks by Monte Carlo simulation. For
each method, 10,000 sets of estimates were drawn from the estimated
sampling distributions (logit scale, t distribution with the method's
degrees of freedom), and we recorded each subgroup's probability of
being among the worst-affected, its probability of ranking first and its
95\% rank interval. Draws were independent across subgroups; this is
exact for regions, whose sampling strata are nested within region, and
approximate for wealth quintiles. The median agreement between the CCA
point ranking and rankings re-drawn from the CCA sampling distribution
provided a sampling-noise benchmark against which method-induced
reordering was judged.

\subsection{Simulation study (RQ3)}\label{simulation-study-rq3}

The simulation study was planned and is reported following the ADEMP
framework (aims, data-generating mechanisms, estimands, methods,
performance measures) \cite{ref31}. The full protocol, including a
documented amendment made after the pilot, is available with the
analysis code.

\textbf{Aims.} To estimate the bias, precision and CI coverage of each
method for survey-weighted stunting and wasting prevalence under MCAR,
MAR and MNAR outcome missingness, and the recovery of the true regional
ranking.

\textbf{Data-generating mechanisms.} A pseudo-population was formed from
the children with valid HAZ and WHZ and observed maternal covariates (n
= 6,247), and its survey-weighted prevalences were taken as the true
values. Each replicate proceeded in two steps. First, a new survey
sample was drawn by the Rao--Wu rescaling bootstrap \cite{ref32}: within
each stratum, n\_h $-$ 1 PSUs were drawn with replacement from the n\_h
available and weights were rescaled, so that replicate-to-replicate
variation reflects the complex sampling design as well as missingness.
One stratum with only two PSUs was merged with the rural stratum of the
same region for this purpose. Second, HAZ and WHZ were deleted jointly,
mirroring the dominant pattern in the survey, where all 178 children
with missing (non-flagged) WHZ had not been measured and also lacked
HAZ. The probability of missingness followed a logistic model under four
mechanisms:

\begin{itemize}
\tightlist
\item
  \textbf{MCAR:} constant probability.
\item
  \textbf{MAR:} the linear predictor of a logistic regression of
  observed missingness on region, residence, wealth quintile, sex, age
  band, mother's interview status and household size, fitted to the
  6,682 children, calibrating the mechanism to the survey's actual
  missingness pattern.
\item
  \textbf{Strong MAR:} log(2) multiplied by each of the standardised
  covariate-predicted logits of stunting and of wasting, from
  pseudo-population logistic regressions on the auxiliary variables used
  by the IPW and MI models. The odds of missingness therefore doubled
  per standard deviation of predicted undernutrition risk.
\item
  \textbf{MNAR:} the MAR linear predictor plus log(2) multiplied by each
  of the negative standardised HAZ and WHZ, so that the odds of
  missingness doubled per standard deviation lower z-score, given
  covariates.
\end{itemize}

The intercept was solved numerically in each replicate to give expected
missingness proportions of 10\%, 30\% and 50\%, yielding 12 scenarios.
The strong-MAR mechanism was added after a pilot (20 replicates per
scenario) showed that the calibrated MAR mechanism induced negligible
CCA bias; it has the same structure and strength as the MNAR mechanism,
driven by predicted risk and not by the outcome itself.

\textbf{Estimands.} National stunting and wasting prevalence (primary);
prevalence by region and wealth quintile, and the true regional ranking
and top-five set (secondary).

\textbf{Methods.} CCA, IPW, MI-PMM and MI-RF as described above, except
that MI used 10 imputations and 10 iterations, and MI-RF used 10 trees
(the mice default), to make the study computationally feasible. Because
no covariates are missing in the pseudo-population, CCA with complete
covariates coincides with CCA and was omitted. The estimate from each
bootstrap sample before deletion (``complete data'') served as a
benchmark for the calibration of the resampling and variance estimation.
In a sub-study (MAR, 30\% missing; 200 replicates), MI-RF was repeated
with 50 trees.

\textbf{Performance measures.} Bias, relative bias, empirical and
average model-based standard errors, their ratio, root mean squared
error, 95\% CI coverage and mean CI width, each with Monte Carlo
standard errors \cite{ref31}; and, for regions, mean absolute bias and
coverage across subgroups, the mean Spearman correlation between
estimated and true rankings, mean top-five overlap and the proportion of
replicates recovering the true top five and the true worst-affected
region. With 500 replicates per scenario, the Monte Carlo standard error
of coverage is 0.97 percentage points at a true coverage of 95\% and at
most 2.2 percentage points. Each replicate used its own random-number
seed and was stored on completion, so that any replicate can be
reproduced individually.

\subsection{Software and
reproducibility}\label{software-and-reproducibility}

Analyses were conducted in R 4.5.1 \cite{ref33} using survey 4.5
\cite{ref22}, mice 3.19.0 \cite{ref13}, ranger 0.18.0 \cite{ref27},
haven 2.5.5, dplyr 1.2.1, ggplot2 4.0.2 and gt 1.3.0; package versions
are recorded with renv. The global random seed was 20210301. The
complete pipeline, from reading the recode files to producing every
table and figure, is scripted; apart from the published reference values
used for validation, no number was entered by hand. Analysis code is
available at \textbf{{[}CONFIRM repository URL{]}}; the DHS microdata
are available from The DHS Program on registration and are not
redistributed.

\section{Results}\label{results}

\subsection{Study population and missing
data}\label{study-population-and-missing-data}

The analysis included 6,682 de facto children aged 0--59 months (Table
1). A valid height-for-age z-score was unavailable for 280 children
(weighted 4.2\%): 257 were not measured and 23 had a flagged value. A
valid weight-for-height z-score was unavailable for 204 children
(weighted 3.1\%): 178 not measured and 26 flagged. All 178 children
without a WHZ measurement had not been measured at all (not present,
140; refused, 18; other, 17; result not recorded, 3) and therefore also
lacked HAZ. Both z-scores were valid for 6,378 children. Among the
auxiliary variables, only mother's age (223 children, 3.3\%) and
mother's education (6 children) had genuinely missing values.

Missingness was not completely at random. The proportion of children
without a valid HAZ was strongly associated with mother's interview
status (Rao--Scott p \textless{} 0.001). It was 2.4\% when the mother
had been interviewed, 9.9\% when the mother did not live in the
household, and 44.6\% when the mother's interview was incomplete. HAZ
missingness also varied by region (p \textless{} 0.001), from 1.1\% in
Ihorombe to 9.9\% in Androy and 9.8\% in Sava. It varied by residence
(urban 6.3\%, rural 3.9\%; p = 0.009) and by age (highest at 48--59
months, 6.9\%; p \textless{} 0.001), and was higher when the mother had
no education (6.3\%) than with primary or secondary education (2.9\%
each; p \textless{} 0.001). It was not associated with wealth quintile
(p = 0.62) or sex (p = 0.82). WHZ missingness was likewise associated
with mother's interview status (p \textless{} 0.001) and residence (p =
0.002), but not with region (p = 0.19).

\subsection{Prevalence by missing-data method
(RQ1)}\label{prevalence-by-missing-data-method-rq1}

Complete-case estimates of national prevalence were 39.8\% (95\% CI
38.0--41.7) for stunting, 12.7\% (11.6--14.0) for severe stunting, 7.7\%
(6.9--8.6) for wasting and 1.5\% (1.2--1.9) for severe wasting, matching
the published survey figures (Table 2). National estimates barely
changed with the missing-data method. The largest difference from CCA
for any method was 0.12 percentage points (pp), for stunting with CCA
restricted to complete covariates (39.7\%). Estimates under IPW, MI-PMM
and MI-RF were within 0.03 pp of CCA for all four indicators. Confidence
intervals were almost identical across methods. The estimated fraction
of missing information in the MI analyses was 0.02--0.05 for national
estimates and at most 0.12 for any subgroup estimate.

Subgroup estimates differed more, but differences remained small
relative to their sampling uncertainty (Table 2). Across the 56 regional
and wealth-quintile estimates of stunting and wasting, the mean absolute
difference from CCA was 0.04--0.27 pp, depending on method and outcome.
All MI-PMM estimates were within 0.5 pp of CCA, as were 96\% of IPW and
MI-RF estimates and 89\% of estimates from CCA with complete covariates.
The largest differences all concerned regional stunting:

\begin{itemize}
\tightlist
\item
  IPW lowered stunting prevalence in Androy from 45.3\% to 43.7\% ($-$1.6
  pp). Androy had the highest proportion of children without a valid HAZ
  (47 of 476).
\item
  MI-RF raised stunting prevalence in Sava from 31.0\% to 32.0\% (+1.1
  pp).
\end{itemize}

All other differences were below 0.9 pp.

Diagnostics did not indicate unstable adjustment. The largest IPW
adjustment factor, the inverse of the estimated response probability,
was 7.1 for HAZ and 3.2 for WHZ. IPW reduced the Kish effective sample
size of respondents only from 4,727 to 4,569 (HAZ) and from 4,785 to
4,728 (WHZ). Imputed values resembled the observed distribution. Among
children with missing HAZ, 41.0\% (MI-PMM) and 40.1\% (MI-RF) of imputed
values were below $-$2, compared with 39.3\% of observed values
(unweighted). For children with flagged HAZ, the figures were 38.9\% and
41.7\%. For WHZ, 8.4\% (MI-PMM) and 7.1\% (MI-RF) of imputed values for
unmeasured children were below $-$2, against 7.8\% of observed values.
Chains for both MI methods mixed well (Additional file 1
\textbf{{[}CONFIRM numbering{]}}).

Excluding flagged children from the target population left conclusions
unchanged (Table S2). National estimates changed by about 0.05 pp or
less for any method, and subgroup estimates by at most 0.8 pp.

\subsection{Equity ordering of regions and wealth quintiles
(RQ2)}\label{equity-ordering-of-regions-and-wealth-quintiles-rq2}

Rankings of regions and wealth quintiles were highly concordant across
methods (Figure 1). For regional stunting, Spearman's $\rho$ between any two
methods was at least 0.993 and Kendall's $\tau$-b at least 0.953, and no
region moved by more than two ranks. For regional wasting, $\rho$ was at
least 0.987 and $\tau$-b at least 0.937, with a largest shift of three ranks.
The ranking of wealth quintiles was identical under all methods for both
outcomes.

The five worst-affected regions were the same under CCA, IPW, MI-PMM and
MI-RF for both outcomes; the top-five set changed only under CCA with
complete covariates.

\begin{itemize}
\tightlist
\item
  \textbf{Stunting.} Analamanga (excluding Antananarivo) replaced Atsimo
  Atsinanana in fifth place under CCA with complete covariates. The two
  regions differed by 0.02 pp under that method (48.28\% vs 48.26\%).
\item
  \textbf{Wasting.} Under CCA with complete covariates, Vatovavy
  Fitovinany fell from fifth to eighth (10.2\% to 9.5\%) and Anosy
  entered the top five.
\item
  \textbf{First place.} The worst-affected region for stunting differed
  under one method only: IPW ranked Itasy (51.57\%) above Vakinankaratra
  (51.56\%). Androy ranked first for wasting under every method.
\end{itemize}

Method-induced reordering was small compared with the uncertainty that
sampling error alone places on the rankings. When CCA estimates were
re-drawn from their sampling distributions, the median agreement with
the CCA point ranking was $\rho$ = 0.93 for stunting and 0.88 for wasting.
The median overlap of the top five was three of five regions for
stunting and four of five for wasting. By contrast, agreement between
methods was $\rho$ $\ge$ 0.987 with a top-five overlap of at least four of five.

No region's position among the five most stunted was secure under CCA.
The highest probability of belonging to the top five was 0.82
(Vakinankaratra), and the probability of ranking first was 0.31 for
Vakinankaratra and 0.30 for Itasy. The 95\% rank intervals of seven
regions included first place. For wasting, only Androy was clearly among
the worst-affected: its probability of being in the top five was 0.97
and of ranking first 0.58, with a 95\% rank interval of 1--6. The 95\%
rank intervals of five regions included first place. The regions whose
worst-affected status depended on the method were among the least
certain. Their probabilities of top-five membership under CCA were
0.53--0.57 for stunting and 0.30--0.36 for wasting.

Among wealth quintiles, stunting prevalence under CCA was highest in the
second-poorest quintile, 43.4\% (40.3--46.6), not the poorest, 42.8\%
(39.8--45.9). The probability that the second-poorest quintile truly
ranked first was 0.57. The richest quintile was clearly the least
affected, 30.0\% (25.7--34.7), and ranked last in at least 97.5\% of
simulated rankings. Wasting prevalence was highest in the poorest
quintile, 9.9\% (8.2--11.9), with a probability of ranking first of
0.90; the gradient across the other quintiles was not monotone. These
orderings were unaffected by the choice of missing-data method.

The flagged-excluded sensitivity analysis gave the same concordance
measures to within one rank shift and identified the same
method-dependent regions (Figure S1).

\subsection{Simulation study (RQ3)}\label{simulation-study-rq3-1}

All 6,200 replicates completed without error, and the imposed
missingness proportions matched their targets in every scenario (mean
achieved 10.0\%, 30.0\% and 50.0\% unweighted). The complete-data
benchmark, the estimate from each bootstrap sample before deletion,
behaved as intended throughout: absolute bias never exceeded 0.07
percentage points (pp), the ratio of model-based to empirical standard
error lay between 0.92 and 1.05, and coverage of nominal 95\% intervals
ranged from 92.8\% to 96.2\%. Departures from nominal coverage below
therefore reflect the missing-data methods, not the resampling design
(Table 3, Table S3; Figure 2).

\textbf{MCAR and calibrated MAR.} When missingness was completely at
random, or followed the mechanism calibrated to the survey's own
missingness pattern, every method was essentially unbiased at all three
rates. For national stunting, absolute bias was at most 0.14 pp under
MCAR and 0.10 pp under MAR (Monte Carlo SE 0.05 pp), and coverage stayed
within 92.2--97.0\% under MCAR. The exception was MI-RF at 50\% MAR
missingness, where standard errors were too small (model-to-empirical SE
ratio 0.85) and coverage fell to 88.6\%. These scenarios reproduce the
near-identical estimates seen on the observed data, where about 4\% of
outcomes were missing.

\textbf{Strong MAR.} When missingness depended strongly on
covariate-predicted undernutrition risk, the methods separated sharply
(Figure 2, Table 3). For national stunting:

\begin{longtable}[]{@{}
  >{\raggedright\arraybackslash}p{(\linewidth - 8\tabcolsep) * \real{0.2000}}
  >{\raggedright\arraybackslash}p{(\linewidth - 8\tabcolsep) * \real{0.2000}}
  >{\raggedright\arraybackslash}p{(\linewidth - 8\tabcolsep) * \real{0.2000}}
  >{\raggedright\arraybackslash}p{(\linewidth - 8\tabcolsep) * \real{0.2000}}
  >{\raggedright\arraybackslash}p{(\linewidth - 8\tabcolsep) * \real{0.2000}}@{}}
\toprule\noalign{}
\begin{minipage}[b]{\linewidth}\raggedright
Missingness
\end{minipage} & \begin{minipage}[b]{\linewidth}\raggedright
CCA
\end{minipage} & \begin{minipage}[b]{\linewidth}\raggedright
IPW
\end{minipage} & \begin{minipage}[b]{\linewidth}\raggedright
MI-PMM
\end{minipage} & \begin{minipage}[b]{\linewidth}\raggedright
MI-RF
\end{minipage} \\
\midrule\noalign{}
\endhead
\bottomrule\noalign{}
\endlastfoot
10\% & $-$0.73 pp (86.4\%) & 0.00 pp (95.6\%) & $-$0.16 pp (95.0\%) & $-$0.24
pp (94.0\%) \\
30\% & $-$2.18 pp (44.2\%) & $-$0.01 pp (96.0\%) & $-$0.52 pp (91.8\%) & $-$0.86
pp (84.2\%) \\
50\% & $-$3.68 pp (11.0\%) & $-$0.04 pp (94.8\%) & $-$0.90 pp (86.4\%) & $-$1.69
pp (63.4\%) \\
\end{longtable}

\emph{Bias in percentage points, with coverage of nominal 95\% CIs in
parentheses.}

IPW removed essentially all of the bias and maintained nominal coverage
at every rate, at the cost of a modest increase in variance (empirical
SE 1.26 pp versus 1.17 pp for CCA at 50\%). Complete-case analysis was
substantially biased downwards and its intervals were badly
miscalibrated: at 50\% missingness its 95\% CI covered the true
prevalence in 11\% of replicates, because its standard errors were
correct for the wrong estimand (SE ratio 1.01). Both imputation methods
removed most but not all of the bias, MI-PMM more successfully than
MI-RF. Root mean squared error at 50\% missingness was 1.26 pp for IPW,
1.52 pp for MI-PMM, 2.05 pp for MI-RF and 3.87 pp for CCA. The same
ordering held for wasting, where CCA bias reached $-$1.11 pp (coverage
54.2\%) at 50\% missingness while IPW remained unbiased ($-$0.04 pp,
coverage 95.2\%).

\textbf{MNAR.} When missingness depended on the child's own z-score
given covariates, no method recovered the truth, as expected: all four
rely on the missing-at-random assumption. For national stunting, bias at
10\%, 30\% and 50\% missingness was $-$2.31, $-$6.25 and $-$10.30 pp for CCA
and $-$1.99, $-$5.69 and $-$9.61 pp for MI-RF, with the other methods in
between; coverage was 33.6--45.2\% at 10\% missingness and 0\% at both
higher rates. Differences between methods were small relative to the
bias they all shared: IPW and MI reduced CCA's bias by at most 8\%. The
corresponding figures for wasting were $-$1.19 to $-$3.52 pp for CCA, with
coverage of 25.4\% at 10\% missingness and 0\% thereafter.

\textbf{Variance estimation under multiple imputation.} MI-RF
understated uncertainty whenever the fraction of missing information was
substantial: its SE ratio fell to 0.81--0.87 in the 50\% scenarios,
against 0.87--0.94 for MI-PMM and 0.93--1.03 for IPW. In the sub-study
(MAR, 30\% missingness, 200 replicates), increasing the number of trees
from 10 to 50 left bias unchanged ($-$0.10 pp versus $-$0.09 pp for
stunting) but improved calibration: the SE ratio rose from 0.93 to 0.99
and coverage from 93.4\% to 96.5\% (Table S5). The shortfall in MI-RF's
variance therefore reflects too few trees rather than the random-forest
imputer as such, but the bias it leaves under strong MAR does not.

\textbf{Subgroup estimates.} Regional estimates showed the same ordering
with larger errors (Table S3). At 50\% missingness under strong MAR,
mean absolute bias across the 23 regions was 0.19 pp for IPW, 1.58 pp
for MI-PMM, 2.37 pp for MI-RF and 2.76 pp for CCA, and mean coverage was
93.6\%, 93.6\%, 90.8\% and 89.7\% respectively. Under MCAR and
calibrated MAR, where CCA is unbiased, imputation was noticeably noisier
at regional level: mean absolute bias was 0.19 pp for CCA against 1.14
pp for MI-PMM and 1.95 pp for MI-RF under MCAR at 50\% missingness,
although coverage remained near nominal for all methods. Imputation
therefore bought robustness to selective missingness at the price of
precision in small domains.

\textbf{Recovery of the equity ranking.} Recovering the true ordering of
the 23 regions was difficult even with complete data: the mean Spearman
correlation between estimated and true regional rankings was 0.92, the
mean overlap of the estimated and true five worst-affected regions was
0.69 (3.5 of 5 regions), the true top-five set was recovered exactly in
5--8\% of replicates, and the truly worst-affected region was identified
in 28--37\%. Missingness and the choice of method changed this little:
at 30\% missingness the mean correlation was 0.90--0.92 in every
mechanism and method, and even at 50\% MNAR missingness it fell only to
0.84--0.87. The dominant obstacle to correct targeting was therefore
sampling variability in a survey designed for 23 domains, not the
treatment of missing data (Table S4).

\section{Discussion}\label{discussion}

\subsection{Principal findings}\label{principal-findings}

In the 2021 Madagascar DHS, the choice of missing-data method barely
affected estimates of child undernutrition. National stunting and
wasting prevalence under inverse probability weighting and two forms of
multiple imputation differed from complete-case estimates by a few
hundredths of a percentage point, and CIs were practically identical.
Subgroup estimates moved more, by up to 1.6 percentage points for
regional stunting, but concordance between regional and wealth-quintile
rankings was very high. The five worst-affected regions were identical
under CCA, IPW, MI-PMM and MI-RF. The two changes to that set occurred
only when analysis was restricted to children with complete maternal
covariates, an approach that changes the target population, not the
estimator.

The more consequential finding concerns uncertainty, not bias. Sampling
error alone reordered regional rankings far more than any change of
method: the median Spearman correlation between the complete-case
ranking and rankings re-drawn from its own sampling distribution was
0.93 for stunting and 0.88 for wasting, whereas agreement between
methods never fell below 0.987. No region's place among the five most
stunted was secure, and the regions whose worst-affected status depended
on the method were precisely those with the most uncertain rank. For
decisions about geographic targeting, the uncertainty inherent in
ranking 23 survey domains outweighed the choice of missing-data method
in this survey.

The simulation showed when that agreement can be expected and when it
cannot. Under MCAR and under the mechanism calibrated to the survey's
own missingness, all four methods were essentially unbiased at 10-50\%
missingness, reproducing what we observed empirically. Under strong MAR,
where missingness depended on covariate-predicted undernutrition risk,
the methods diverged sharply: at 50\% missingness complete-case analysis
under-estimated stunting by 3.68 pp with 11\% CI coverage, IPW was
unbiased (-0.04 pp) with 94.8\% coverage, and multiple imputation
removed most but not all of the bias (-0.90 pp for MI-PMM, -1.69 pp for
MI-RF). Under MNAR no method recovered the truth: all were biased by
about -10 pp at 50\% missingness with zero coverage. Recovery of the
regional ranking, by contrast, was poor in every scenario, including
with complete data.

\subsection{Why the methods agreed}\label{why-the-methods-agreed}

Three features of these data explain the agreement, and they indicate
when it cannot be assumed. First, outcome missingness was low: about 4\%
of children lacked a valid HAZ and 3\% a valid WHZ. The bias of a
complete-case estimate is bounded by the proportion missing, so with so
few missing outcomes even a selective mechanism can shift national
prevalence only slightly. Second, although missingness was clearly not
completely at random, it was concentrated in small groups. It reached
44.6\% among children whose mother's interview was incomplete, but that
group comprised 193 children, whereas missingness was 2.4\% among the
large majority whose mother was interviewed. Such concentrated selection
has little leverage on the national estimate unless those groups differ
greatly in outcome. Third, IPW and both MI methods drew on largely the
same auxiliary information, so under MAR given those covariates they
target the same estimand and would be expected to agree with each other.
They differ from CCA only to the extent that the covariates predicting
missingness also predict the outcome.

The largest subgroup differences arose where these conditions weakened.
Androy had the highest proportion of children without a valid HAZ, and
IPW lowered its stunting estimate by 1.6 percentage points. That
region-specific correction is plausible, but it rests entirely on the
MAR assumption within a single domain of fewer than 500 children, and
its magnitude was well within the region's sampling uncertainty. It
illustrates that national agreement does not guarantee agreement within
subgroups where missingness is concentrated.

These findings agree with guidance that the proportion of missing data
alone is a poor guide to analysis strategy, and that complete-case
analysis can be adequate when missingness is low and plausibly unrelated
to the outcome given the design \cite{ref34,ref35}. They also show that
such adequacy should be demonstrated, not assumed, particularly for
subnational estimates. A framework such as TARMOS, which asks analysts
to state the plausible missingness mechanisms and check the robustness
of conclusions, fits this purpose well \cite{ref16}.

The simulation supports this reading. The mechanism calibrated to the
survey's own missingness produced almost no bias in complete-case
estimates (at most 0.09 pp for stunting at 50\% missingness), which is
why the methods agreed on the observed data. Selective missingness only
became consequential when it was tied to the outcome's own predictors:
the strong-MAR mechanism, with the same structure and strength but
driven by predicted undernutrition risk, produced complete-case bias of
-2.18 pp at 30\% missingness. The practical implication is that the
observed agreement is a property of this survey's missingness pattern,
not a general result.

\subsection{Flagged measurements}\label{flagged-measurements}

Flagged z-scores are handled inconsistently in practice. DHS reports
exclude them, and cleaning criteria are known to change malnutrition
prevalence estimates \cite{ref10}. We treated them as a distinct form of
missingness in the main analysis and excluded them in a sensitivity
analysis. With 23 flagged HAZ and 26 flagged WHZ values among 6,682
children, the choice made no material difference here: national
estimates changed by about 0.05 percentage points or less. In surveys
with poorer anthropometric data quality, where flag rates are higher
\cite{ref8,ref9}, the choice could matter more. Imputing flagged values
under MAR assumes that the plausible range of the true measurement can
be predicted from covariates. That assumption is questionable if flags
arise from extreme true values, and less so if they arise from recording
error.

\subsection{Implications for survey analysis and
targeting}\label{implications-for-survey-analysis-and-targeting}

Our results suggest three practical recommendations for analysts using
DHS-type surveys to rank subnational units.

\begin{enumerate}
\def\labelenumi{\arabic{enumi}.}
\item
  \textbf{Report rank uncertainty.} Point rankings of survey domains
  should be accompanied by probabilities of membership in the targeted
  set or by rank intervals, as has long been recommended for league
  tables of institutional performance \cite{ref18,ref19}. In this
  survey, even the region with the highest stunting prevalence had only
  about a one-in-three probability of truly ranking first.
\item
  \textbf{Check robustness at the subgroup level.} Missing-data
  sensitivity analyses should be reported for the subgroup estimates
  used for targeting, not only for national prevalence. Regions with the
  highest missingness are those where methods are most likely to
  diverge.
\item
  \textbf{Avoid restricting to complete covariates.} Restricting to
  records complete for auxiliary variables changed the worst-affected
  set in both outcomes. Here it excluded children whose mother lived in
  the household but could not be linked to an interview, a group with
  very high outcome missingness, and it changes the population to which
  estimates apply. Structurally inapplicable values, such as maternal
  characteristics for children whose mother lives elsewhere, should be
  coded as such, not treated as missing.
\item
  \textbf{Prefer response weighting, or imputation with well-calibrated
  variance, when missingness is selective.} In the simulation IPW was
  the only method that stayed unbiased with nominal coverage under
  strong MAR, while both imputation approaches retained residual bias.
  MI-RF additionally understated uncertainty: its standard errors were
  13-19\% too small in the highest-missingness scenarios, and coverage
  fell to 63\% under strong MAR at 50\% missingness. Where random-forest
  imputation is used, the number of trees should be checked: raising it
  from 10 (the mice default) to 50 restored calibration in our
  sub-study, though it did not remove the bias.
\end{enumerate}

Where ranking precision is itself the problem, borrowing strength across
areas through small-area estimation methods may be more useful than
refining the treatment of a few percent of missing outcomes
\cite{ref36}.

\subsection{Strengths and limitations}\label{strengths-and-limitations}

The analysis has several strengths:

\begin{itemize}
\tightlist
\item
  \textbf{Validated baseline.} Before any method was applied, the
  complete-case baseline reproduced every cell of the survey's published
  nutrition table and its sampling errors.
\item
  \textbf{Design-based inference throughout,} including within multiple
  imputation.
\item
  \textbf{Missingness taxonomy.} Structural missingness (children whose
  mother does not live in the household) was distinguished from genuine
  missingness, and flagged values from unmeasured children.
\item
  \textbf{Transparent simulation.} The simulation followed a
  pre-specified protocol that recorded one amendment made after the
  pilot.
\item
  \textbf{Reproducibility.} The whole pipeline is scripted and openly
  available; no microdata are redistributed.
\end{itemize}

The study also has limitations:

\begin{itemize}
\tightlist
\item
  \textbf{Untestable MAR assumption.} IPW and MI assume MAR given the
  auxiliary variables. The observed data cannot rule out MNAR, for
  example if children who were absent at measurement were systematically
  more or less undernourished than similar children who were measured.
  The observed-data analyses therefore compare methods that share this
  assumption; they do not establish that any is unbiased. Sensitivity
  analyses that relax MAR, such as delta-adjusted multiple imputation
  \cite{ref37} or selection models \cite{ref17}, would complement this
  work. The MNAR simulation quantifies what is at stake rather than
  resolving it: when missingness depended on the z-score itself, every
  method was biased by a similar amount and no confidence interval
  covered the truth at 30\% missingness or above, with IPW and MI
  reducing complete-case bias by at most 8\%.
\item
  \textbf{Imputation model.} The model included stratum-defining
  variables and the design weight but not cluster-level random effects,
  which may slightly understate between-cluster variation in imputed
  values. Weight-for-age was excluded to achieve convergence.
\item
  \textbf{Variance estimation.} IPW treated estimated weights as fixed.
  MI used 20 imputations, and we did not quantify the Monte Carlo error
  of MI point estimates, so differences of a few hundredths of a
  percentage point between methods should not be interpreted.
\item
  \textbf{Ranking analysis.} The top-five threshold is a convenient
  illustration, not a policy rule, and simulated rankings of wealth
  quintiles assumed independence across quintiles.
\item
  \textbf{Survey scope.} The survey's 23 domains predate the division of
  Vatovavy Fitovinany into two regions. The findings come from a single
  survey with low missingness and may not transfer to surveys with
  higher or differently structured missingness. The simulation extends
  the range of missingness studied to 50\%, but within the same
  pseudo-population and design, and it deletes HAZ and WHZ together; the
  smaller HAZ-only pattern seen in the survey (1.2\% of children) was
  not simulated.
\end{itemize}

\subsection{Conclusions}\label{conclusions}

In the 2021 Madagascar DHS, complete-case analysis, inverse probability
weighting and multiple imputation gave practically identical national
estimates of stunting and wasting and nearly identical subnational
rankings. The worst-affected set changed only under restriction to
complete covariates. Sampling uncertainty in subnational rankings was
substantially larger than the effect of the missing-data method, and
should be reported alongside any ranking used for targeting. In
simulation, this equivalence held under missingness completely at random
and under the mechanism calibrated to the survey, but not when
missingness was strongly tied to covariate-predicted undernutrition
risk: there only inverse probability weighting remained unbiased with
valid coverage, while complete-case analysis was substantially biased.
No method compensated for missingness that depended on the measurement
itself.

\section{Declarations (excerpt)}\label{declarations-excerpt}

\textbf{Ethics approval and consent to participate.} This study is a
secondary analysis of anonymised survey data. The EDSMD-V 2021 protocol,
including biomarker collection, was reviewed and approved by the
Biomedical Research Ethics Committee of the Ministry of Public Health of
Madagascar and by the ICF ethics committee, and informed consent was
obtained before biomarker measurement \cite{ref6}. \textbf{{[}CONFIRM
wording against DHS standard consent procedures and journal
requirements.{]}} Data were obtained from The DHS Program under its
terms of use.

\clearpage
\begin{thebibliography}{99}
\bibitem{ref1} Black RE, Victora CG, Walker SP, Bhutta ZA, Christian P, de Onis M, et al. Maternal and child undernutrition and overweight in low-income and middle-income countries. Lancet. 2013;382(9890):427–51.
\bibitem{ref2} de Onis M, Branca F. Childhood stunting: a global perspective. Matern Child Nutr. 2016;12(Suppl 1):12–26.
\bibitem{ref3} UNICEF, World Health Organization, World Bank Group. Levels and trends in child malnutrition: UNICEF/WHO/World Bank Group joint child malnutrition estimates, key findings of the 2023 edition. Geneva: World Health Organization; 2023.
\bibitem{ref4} Victora CG, Christian P, Vidaletti LP, Gatica-Domínguez G, Menon P, Black RE. Revisiting maternal and child undernutrition in low-income and middle-income countries: variable progress towards an unfinished agenda. Lancet. 2021;397(10282):1388–99.
\bibitem{ref5} World Health Organization. Global nutrition targets 2025: stunting policy brief. Geneva: World Health Organization; 2014.
\bibitem{ref6} Institut National de la Statistique (INSTAT), ICF. Enquête Démographique et de Santé à Madagascar, 2021. Antananarivo, Madagascar and Rockville, Maryland, USA: INSTAT and ICF; 2022.
\bibitem{ref7} Corsi DJ, Neuman M, Finlay JE, Subramanian SV. Demographic and health surveys: a profile. Int J Epidemiol. 2012;41(6):1602–13.
\bibitem{ref8} Assaf S, Kothari MT, Pullum T. An assessment of the quality of DHS anthropometric data, 2005–2014. DHS Methodological Reports No. 16. Rockville, Maryland, USA: ICF International; 2015.
\bibitem{ref9} Perumal N, Namaste S, Qamar H, Aimone A, Bassani DG, Roth DE. Anthropometric data quality assessment in multisurvey studies of child growth. Am J Clin Nutr. 2020;112(Suppl 2):806S–815S.
\bibitem{ref10} Crowe S, Seal A, Grijalva-Eternod C, Kerac M. Effect of nutrition survey 'cleaning criteria' on estimates of malnutrition prevalence and disease burden: secondary data analysis. PeerJ. 2014;2:e380.
\bibitem{ref11} Little RJA, Rubin DB. Statistical analysis with missing data. 3rd ed. Hoboken, NJ: Wiley; 2019.
\bibitem{ref12} Seaman SR, White IR. Review of inverse probability weighting for dealing with missing data. Stat Methods Med Res. 2013;22(3):278–95.
\bibitem{ref13} van Buuren S, Groothuis-Oudshoorn K. mice: Multivariate Imputation by Chained Equations in R. J Stat Softw. 2011;45(3):1–67.
\bibitem{ref14} Doove LL, Van Buuren S, Dusseldorp E. Recursive partitioning for missing data imputation in the presence of interaction effects. Comput Stat Data Anal. 2014;72:92–104.
\bibitem{ref15} Shah AD, Bartlett JW, Carpenter J, Nicholas O, Hemingway H. Comparison of random forest and parametric imputation models for imputing missing data using MICE: a CALIBER study. Am J Epidemiol. 2014;179(6):764–74.
\bibitem{ref16} Lee KJ, Tilling KM, Cornish RP, Little RJA, Bell ML, Goetghebeur E, et al. Framework for the treatment and reporting of missing data in observational studies: the Treatment And Reporting of Missing data in Observational Studies framework. J Clin Epidemiol. 2021;134:79–88.
\bibitem{ref17} Bärnighausen T, Bor J, Wandira-Kazibwe S, Canning D. Correcting HIV prevalence estimates for survey nonparticipation using Heckman-type selection models. Epidemiology. 2011;22(1):27–35.
\bibitem{ref18} Goldstein H, Spiegelhalter DJ. League tables and their limitations: statistical issues in comparisons of institutional performance. J R Stat Soc Ser A Stat Soc. 1996;159(3):385–443.
\bibitem{ref19} Marshall EC, Spiegelhalter DJ. Reliability of league tables of in vitro fertilisation clinics: retrospective analysis of live birth rates. BMJ. 1998;316(7146):1701–5.
\bibitem{ref20} Croft TN, Allen CK, Zachary BW, et al. Guide to DHS Statistics: DHS-8. Rockville, Maryland, USA: ICF; 2023.
\bibitem{ref21} WHO Multicentre Growth Reference Study Group. WHO Child Growth Standards based on length/height, weight and age. Acta Paediatr Suppl. 2006;450:76–85.
\bibitem{ref22} Lumley T. Analysis of complex survey samples. J Stat Softw. 2004;9(8):1–19.
\bibitem{ref23} Rao JNK, Scott AJ. On chi-squared tests for multiway contingency tables with cell proportions estimated from survey data. Ann Stat. 1984;12(1):46–60.
\bibitem{ref24} White IR, Royston P, Wood AM. Multiple imputation using chained equations: issues and guidance for practice. Stat Med. 2011;30(4):377–99.
\bibitem{ref25} Morris TP, White IR, Royston P. Tuning multiple imputation by predictive mean matching and local residual draws. BMC Med Res Methodol. 2014;14:75.
\bibitem{ref26} Reiter JP, Raghunathan TE, Kinney SK. The importance of modeling the sampling design in multiple imputation for missing data. Surv Methodol. 2006;32(2):143–9.
\bibitem{ref27} Wright MN, Ziegler A. ranger: a fast implementation of random forests for high dimensional data in C++ and R. J Stat Softw. 2017;77(1):1–17.
\bibitem{ref28} Rubin DB. Multiple imputation for nonresponse in surveys. New York: Wiley; 1987.
\bibitem{ref29} Marshall A, Altman DG, Holder RL, Royston P. Combining estimates of interest in prognostic modelling studies after multiple imputation: current practice and guidelines. BMC Med Res Methodol. 2009;9:57.
\bibitem{ref30} Barnard J, Rubin DB. Small-sample degrees of freedom with multiple imputation. Biometrika. 1999;86(4):948–55.
\bibitem{ref31} Morris TP, White IR, Crowther MJ. Using simulation studies to evaluate statistical methods. Stat Med. 2019;38(11):2074–102.
\bibitem{ref32} Rao JNK, Wu CFJ. Resampling inference with complex survey data. J Am Stat Assoc. 1988;83(401):231–41.
\bibitem{ref33} R Core Team. R: a language and environment for statistical computing. Vienna, Austria: R Foundation for Statistical Computing; 2025.
\bibitem{ref34} Hughes RA, Heron J, Sterne JAC, Tilling K. Accounting for missing data in statistical analyses: multiple imputation is not always the answer. Int J Epidemiol. 2019;48(4):1294–304.
\bibitem{ref35} Jakobsen JC, Gluud C, Wetterslev J, Winkel P. When and how should multiple imputation be used for handling missing data in randomised clinical trials – a practical guide with flowcharts. BMC Med Res Methodol. 2017;17:162.
\bibitem{ref36} Mercer LD, Wakefield J, Pantazis A, Lutambi AM, Masanja H, Clark S. Space–time smoothing of complex survey data: small area estimation for child mortality. Ann Appl Stat. 2015;9(4):1889–905.
\bibitem{ref37} Leacy FP, Floyd S, Yates TA, White IR. Analyses of sensitivity to the missing-at-random assumption using multiple imputation with delta adjustment: application to a tuberculosis/HIV prevalence survey with incomplete HIV-status data. Am J Epidemiol. 2017;185(4):304–15.
\end{thebibliography}

\clearpage
\section*{Tables}

% --- table1_sample_missingness
\begingroup
\fontsize{8.0pt}{10.0pt}\selectfont
\setlength{\LTpost}{0mm}
\begin{longtable}{l|rrrr}
\caption*{
{\fontsize{20}{25}\selectfont  Table 1. Characteristics of de facto children aged 0–59 months and outcome missingness, Madagascar DHS 2021\fontsize{8}{10}\selectfont }
} \\ 
\toprule
 &  &  & \multicolumn{2}{c}{{No valid z-score, weighted \%}} \\ 
\cmidrule(lr){4-5}
 & n & Weighted \% & HAZ & WHZ \\ 
\midrule\endhead\addlinespace[2.5pt]
\multicolumn{5}{l}{All children} \\[2.5pt] 
\midrule\addlinespace[2.5pt]
 & 6682 & 100.0 & 4.2 & 3.1 \\ 
\midrule\addlinespace[2.5pt]
\multicolumn{5}{l}{Sex} \\[2.5pt] 
\midrule\addlinespace[2.5pt]
Male & 3377 & 50.3 & 4.2 & 3.0 \\ 
Female & 3305 & 49.7 & 4.3 & 3.2 \\ 
\midrule\addlinespace[2.5pt]
\multicolumn{5}{l}{Age (months)} \\[2.5pt] 
\midrule\addlinespace[2.5pt]
<6 & 732 & 10.8 & 3.8 & 5.3 \\ 
6-8 & 332 & 5.0 & 2.7 & 2.5 \\ 
9-11 & 348 & 5.4 & 2.0 & 1.7 \\ 
12-17 & 687 & 10.4 & 3.0 & 2.1 \\ 
18-23 & 592 & 8.8 & 2.4 & 2.4 \\ 
24-35 & 1329 & 19.5 & 4.2 & 2.8 \\ 
36-47 & 1346 & 20.8 & 4.3 & 2.9 \\ 
48-59 & 1316 & 19.5 & 6.9 & 3.8 \\ 
\midrule\addlinespace[2.5pt]
\multicolumn{5}{l}{Residence} \\[2.5pt] 
\midrule\addlinespace[2.5pt]
Urban & 1248 & 15.0 & 6.3 & 5.2 \\ 
Rural & 5434 & 85.0 & 3.9 & 2.7 \\ 
\midrule\addlinespace[2.5pt]
\multicolumn{5}{l}{Region} \\[2.5pt] 
\midrule\addlinespace[2.5pt]
Antananarivo & 319 & 3.5 & 5.7 & 5.6 \\ 
Analamanga (excl. Antananarivo) & 223 & 7.8 & 2.9 & 3.7 \\ 
Vakinankaratra & 274 & 7.2 & 3.8 & 4.5 \\ 
Itasy & 261 & 3.0 & 1.2 & 1.6 \\ 
Bongolava & 325 & 2.7 & 3.0 & 3.0 \\ 
Haute Matsiatra & 259 & 5.0 & 3.2 & 1.0 \\ 
Amoron'i Mania & 287 & 2.8 & 2.4 & 2.8 \\ 
Vatovavy Fitovinany & 289 & 4.7 & 8.8 & 5.3 \\ 
Ihorombe & 322 & 2.0 & 1.1 & 1.1 \\ 
Atsimo Atsinanana & 376 & 4.4 & 1.8 & 1.7 \\ 
Atsinanana & 272 & 5.8 & 3.9 & 2.6 \\ 
Analanjirofo & 294 & 5.6 & 2.1 & 1.7 \\ 
Alaotra Mangoro & 258 & 4.1 & 3.2 & 2.9 \\ 
Boeny & 221 & 3.2 & 1.5 & 1.5 \\ 
Sofia & 225 & 5.8 & 2.7 & 3.1 \\ 
Betsiboka & 271 & 1.5 & 4.1 & 3.7 \\ 
Melaky & 277 & 1.1 & 3.1 & 1.8 \\ 
Atsimo Andrefana & 387 & 9.6 & 6.3 & 4.7 \\ 
Androy & 476 & 5.8 & 9.9 & 3.5 \\ 
Anosy & 333 & 3.9 & 1.8 & 1.3 \\ 
Menabe & 291 & 3.2 & 2.3 & 1.8 \\ 
Diana & 210 & 2.5 & 3.7 & 2.7 \\ 
Sava & 232 & 4.9 & 9.8 & 3.8 \\ 
\midrule\addlinespace[2.5pt]
\multicolumn{5}{l}{Wealth quintile} \\[2.5pt] 
\midrule\addlinespace[2.5pt]
Poorest & 1832 & 25.4 & 4.4 & 2.6 \\ 
Poorer & 1495 & 21.4 & 4.9 & 3.2 \\ 
Middle & 1324 & 21.3 & 3.3 & 2.5 \\ 
Richer & 1024 & 16.8 & 3.8 & 3.3 \\ 
Richest & 1007 & 15.1 & 4.8 & 4.5 \\ 
\midrule\addlinespace[2.5pt]
\multicolumn{5}{l}{Mother's interview status} \\[2.5pt] 
\midrule\addlinespace[2.5pt]
Interviewed & 5890 & 88.2 & 2.4 & 1.7 \\ 
Not de facto & 30 & 0.5 & 6.1 & 6.9 \\ 
Interview incomplete & 193 & 2.9 & 44.6 & 42.4 \\ 
Not in household & 569 & 8.4 & 9.9 & 4.5 \\ 
\midrule\addlinespace[2.5pt]
\multicolumn{5}{l}{Mother's education} \\[2.5pt] 
\midrule\addlinespace[2.5pt]
None & 1574 & 21.3 & 6.3 & 3.2 \\ 
Primary & 2618 & 39.7 & 2.9 & 2.6 \\ 
Secondary+ & 1915 & 30.7 & 2.9 & 3.2 \\ 
Not in household & 569 & 8.4 & 9.9 & 4.5 \\ 
— & 6 & 0.1 & 39.5 & 39.5 \\ 
\bottomrule
\end{longtable}
\begin{minipage}{\linewidth}
\vspace{.05em}
n unweighted; percentages weighted by the DHS sample weight. No valid z-score = not measured or WHO-flagged. HAZ: 257 not measured/missing, 23 flagged. WHZ: 178 not measured/missing, 26 flagged. Mother's age unknown (mother in household, not linked): 223. Mother's education 'don't know': 6.\\
\end{minipage}
\endgroup

\clearpage

% --- table2_prevalence_by_method
\begin{landscape}
\begingroup
\fontsize{8.0pt}{10.0pt}\selectfont
\setlength{\LTpost}{0mm}
\begin{longtable}{l|rlllll}
\caption*{
{\fontsize{20}{25}\selectfont  Table 2. Prevalence of stunting and wasting by missing-data method, Madagascar DHS 2021\fontsize{8}{10}\selectfont }
} \\ 
\toprule
 &  & \multicolumn{5}{c}{{Prevalence, \% (95\% CI)}} \\ 
\cmidrule(lr){3-7}
 & n valid / n & CCA & CCA (outcome + covariates) & IPW & MI-PMM & MI-RF \\ 
\midrule\endhead\addlinespace[2.5pt]
\multicolumn{7}{l}{Stunting — National} \\[2.5pt] 
\midrule\addlinespace[2.5pt]
National & 6402 / 6682 & 39.8 (38.0–41.7) & 39.7 (37.9–41.6) & 39.8 (38.0–41.7) & 39.9 (38.1–41.7) & 39.9 (38.1–41.7) \\ 
\midrule\addlinespace[2.5pt]
\multicolumn{7}{l}{Stunting — Region} \\[2.5pt] 
\midrule\addlinespace[2.5pt]
Antananarivo & 300 / 319 & 44.0 (39.1–49.1) & 43.4 (38.3–48.7) & 44.4 (39.5–49.5) & 44.2 (39.3–49.3) & 43.5 (38.4–48.8) \\ 
Analamanga (excl. Antananarivo) $\dagger$ & 217 / 223 & 48.3 (40.2–56.5) & 48.3 (40.2–56.4) & 48.3 (40.1–56.5) & 48.4 (40.3–56.6) & 48.2 (40.1–56.4) \\ 
Vakinankaratra & 263 / 274 & 51.9 (43.3–60.4) & 52.1 (43.3–60.8) & 51.6 (43.1–59.9) & 51.8 (43.3–60.1) & 51.6 (43.1–60.1) \\ 
Itasy & 258 / 261 & 51.6 (42.7–60.3) & 51.5 (42.8–60.2) & 51.6 (42.7–60.3) & 51.4 (42.6–60.1) & 51.5 (42.7–60.2) \\ 
Bongolava & 314 / 325 & 37.6 (30.6–45.2) & 37.5 (30.4–45.3) & 37.6 (30.6–45.2) & 37.7 (30.7–45.2) & 37.6 (30.6–45.1) \\ 
Haute Matsiatra & 250 / 259 & 49.2 (41.9–56.5) & 49.2 (41.9–56.5) & 49.0 (41.7–56.4) & 49.6 (42.4–56.8) & 49.1 (41.9–56.4) \\ 
Amoron'i Mania & 280 / 287 & 49.1 (41.6–56.6) & 49.0 (41.4–56.7) & 49.5 (42.0–57.0) & 49.1 (41.7–56.7) & 49.1 (41.7–56.6) \\ 
Vatovavy Fitovinany & 261 / 289 & 45.7 (38.5–53.1) & 45.8 (38.4–53.3) & 46.2 (38.7–53.8) & 45.7 (38.7–52.9) & 45.7 (38.7–52.9) \\ 
Ihorombe & 318 / 322 & 32.6 (26.0–40.0) & 32.7 (26.1–40.1) & 32.6 (26.1–40.0) & 32.5 (26.0–39.9) & 32.7 (26.2–39.9) \\ 
Atsimo Atsinanana $\dagger$ & 366 / 376 & 48.5 (43.1–54.0) & 48.3 (42.8–53.8) & 48.5 (43.1–54.0) & 48.4 (43.0–53.9) & 48.4 (42.9–53.9) \\ 
Atsinanana & 259 / 272 & 37.6 (31.1–44.6) & 37.0 (30.5–43.9) & 38.3 (31.4–45.8) & 38.1 (31.7–45.0) & 38.0 (31.5–44.9) \\ 
Analanjirofo & 287 / 294 & 36.2 (28.5–44.6) & 36.2 (28.4–44.8) & 36.3 (28.6–44.7) & 36.2 (28.5–44.6) & 36.2 (28.7–44.5) \\ 
Alaotra Mangoro & 250 / 258 & 44.9 (37.8–52.1) & 45.0 (37.7–52.5) & 44.8 (37.9–51.9) & 44.8 (37.8–52.0) & 44.7 (37.8–51.8) \\ 
Boeny & 217 / 221 & 28.8 (22.1–36.6) & 28.4 (21.8–36.1) & 29.0 (22.2–36.8) & 28.7 (22.0–36.4) & 28.9 (22.2–36.6) \\ 
Sofia & 218 / 225 & 28.1 (20.8–36.9) & 27.3 (20.0–36.1) & 28.4 (21.0–37.1) & 28.1 (20.9–36.7) & 28.3 (21.1–36.8) \\ 
Betsiboka & 261 / 271 & 29.6 (23.8–36.1) & 29.9 (24.1–36.3) & 29.4 (23.4–36.1) & 29.7 (23.9–36.2) & 29.7 (23.8–36.3) \\ 
Melaky & 267 / 277 & 22.3 (19.0–25.9) & 22.0 (18.9–25.6) & 22.5 (19.1–26.2) & 22.5 (19.0–26.4) & 22.3 (19.0–26.1) \\ 
Atsimo Andrefana & 363 / 387 & 26.7 (20.7–33.8) & 26.5 (20.7–33.3) & 27.1 (20.8–34.5) & 26.9 (21.0–33.7) & 27.1 (21.1–34.1) \\ 
Androy & 429 / 476 & 45.3 (40.6–50.1) & 45.7 (40.8–50.8) & 43.7 (38.7–48.9) & 45.2 (40.5–50.0) & 45.3 (40.5–50.2) \\ 
Anosy & 327 / 333 & 46.8 (38.5–55.3) & 46.7 (38.5–55.1) & 46.9 (38.6–55.5) & 46.7 (38.5–55.1) & 46.8 (38.5–55.2) \\ 
Menabe & 284 / 291 & 34.9 (29.6–40.6) & 34.1 (28.6–40.1) & 35.0 (29.7–40.6) & 34.9 (29.7–40.6) & 35.0 (29.7–40.7) \\ 
Diana & 203 / 210 & 21.7 (15.0–30.3) & 21.2 (14.5–29.9) & 21.8 (15.1–30.4) & 21.9 (15.2–30.6) & 21.8 (15.1–30.4) \\ 
Sava & 210 / 232 & 31.0 (24.6–38.1) & 31.2 (24.8–38.3) & 31.2 (24.6–38.7) & 31.2 (24.6–38.7) & 32.0 (25.7–39.1) \\ 
\midrule\addlinespace[2.5pt]
\multicolumn{7}{l}{Stunting — Wealth quintile} \\[2.5pt] 
\midrule\addlinespace[2.5pt]
Poorest & 1749 / 1832 & 42.8 (39.8–45.9) & 42.8 (39.8–46.0) & 42.3 (39.3–45.4) & 42.8 (39.9–45.9) & 42.8 (39.8–45.8) \\ 
Poorer & 1424 / 1495 & 43.4 (40.3–46.6) & 43.2 (40.1–46.4) & 43.7 (40.6–47.0) & 43.3 (40.2–46.5) & 43.5 (40.4–46.7) \\ 
Middle & 1285 / 1324 & 40.1 (36.6–43.8) & 39.8 (36.2–43.4) & 40.3 (36.7–44.1) & 40.2 (36.6–43.9) & 40.1 (36.5–43.8) \\ 
Richer & 986 / 1024 & 39.4 (35.5–43.4) & 39.7 (35.8–43.7) & 39.2 (35.2–43.3) & 39.2 (35.3–43.3) & 39.4 (35.5–43.4) \\ 
Richest & 958 / 1007 & 30.0 (25.7–34.7) & 29.6 (25.3–34.3) & 30.0 (25.8–34.7) & 30.3 (26.0–34.8) & 30.1 (25.9–34.6) \\ 
\midrule\addlinespace[2.5pt]
\multicolumn{7}{l}{Wasting — National} \\[2.5pt] 
\midrule\addlinespace[2.5pt]
National & 6478 / 6682 & 7.7 (6.9–8.6) & 7.7 (6.9–8.6) & 7.7 (6.9–8.6) & 7.7 (6.9–8.6) & 7.7 (6.8–8.6) \\ 
\midrule\addlinespace[2.5pt]
\multicolumn{7}{l}{Wasting — Region} \\[2.5pt] 
\midrule\addlinespace[2.5pt]
Antananarivo & 300 / 319 & 3.7 (2.1–6.5) & 3.8 (2.1–6.6) & 3.6 (2.0–6.3) & 3.7 (2.1–6.6) & 3.7 (2.1–6.6) \\ 
Analamanga (excl. Antananarivo) & 215 / 223 & 5.1 (2.6–9.7) & 4.7 (2.3–9.4) & 5.3 (2.7–10.1) & 5.1 (2.6–9.7) & 5.1 (2.6–9.6) \\ 
Vakinankaratra & 261 / 274 & 4.1 (2.0–8.4) & 4.2 (2.0–8.6) & 4.1 (1.9–8.3) & 4.3 (2.1–8.5) & 4.2 (2.0–8.6) \\ 
Itasy & 257 / 261 & 3.9 (2.2–6.8) & 4.0 (2.3–6.9) & 4.0 (2.3–6.9) & 4.1 (2.3–7.1) & 4.0 (2.3–7.0) \\ 
Bongolava & 314 / 325 & 7.4 (5.0–10.7) & 7.3 (4.8–10.9) & 7.5 (5.1–10.8) & 7.5 (5.1–10.8) & 7.4 (5.0–10.7) \\ 
Haute Matsiatra & 256 / 259 & 2.4 (0.6–9.2) & 2.5 (0.6–9.4) & 2.4 (0.6–9.1) & 2.4 (0.6–9.1) & 2.4 (0.6–9.1) \\ 
Amoron'i Mania & 279 / 287 & 2.6 (1.4–4.9) & 2.7 (1.4–5.0) & 2.6 (1.4–4.8) & 2.7 (1.4–5.0) & 2.6 (1.4–5.0) \\ 
Vatovavy Fitovinany $\dagger$ & 271 / 289 & 10.2 (7.4–13.9) & 9.5 (6.5–13.7) & 10.5 (7.5–14.4) & 10.4 (7.4–14.3) & 10.0 (7.3–13.7) \\ 
Ihorombe & 318 / 322 & 5.2 (2.9–9.1) & 5.2 (2.9–9.2) & 5.2 (2.9–9.1) & 5.2 (2.9–9.1) & 5.2 (2.9–9.1) \\ 
Atsimo Atsinanana & 368 / 376 & 12.4 (9.5–15.9) & 12.5 (9.6–16.1) & 12.3 (9.5–15.8) & 12.3 (9.5–15.9) & 12.4 (9.5–16.0) \\ 
Atsinanana & 264 / 272 & 8.6 (5.1–14.2) & 8.6 (5.0–14.3) & 8.6 (5.2–14.1) & 8.7 (5.2–14.3) & 8.6 (5.1–14.0) \\ 
Analanjirofo & 288 / 294 & 11.3 (7.3–17.0) & 11.1 (7.2–16.8) & 11.4 (7.4–17.2) & 11.2 (7.3–16.8) & 11.1 (7.2–16.7) \\ 
Alaotra Mangoro & 251 / 258 & 3.7 (1.9–7.2) & 3.8 (2.0–7.3) & 3.7 (1.9–7.1) & 3.8 (1.9–7.3) & 3.8 (1.9–7.2) \\ 
Boeny & 217 / 221 & 12.4 (8.3–18.2) & 12.7 (8.5–18.4) & 12.4 (8.3–18.1) & 12.3 (8.2–18.0) & 12.4 (8.3–18.2) \\ 
Sofia & 218 / 225 & 9.8 (6.1–15.2) & 10.0 (6.3–15.4) & 9.6 (6.0–15.0) & 9.7 (6.1–15.0) & 9.7 (6.1–15.1) \\ 
Betsiboka & 262 / 271 & 7.3 (4.2–12.5) & 6.9 (4.2–11.2) & 7.5 (4.2–13.2) & 7.3 (4.2–12.4) & 7.3 (4.2–12.3) \\ 
Melaky & 270 / 277 & 5.9 (3.1–10.9) & 6.0 (3.2–11.1) & 5.8 (3.0–10.9) & 5.9 (3.1–10.9) & 5.9 (3.1–10.9) \\ 
Atsimo Andrefana & 370 / 387 & 6.8 (4.2–11.0) & 6.7 (4.0–11.0) & 7.0 (4.3–11.1) & 6.9 (4.2–10.9) & 6.8 (4.2–10.9) \\ 
Androy & 461 / 476 & 15.2 (10.8–20.8) & 15.6 (11.2–21.3) & 14.9 (10.7–20.5) & 15.1 (10.8–20.7) & 15.1 (10.8–20.7) \\ 
Anosy $\dagger$ & 328 / 333 & 9.9 (7.1–13.7) & 10.1 (7.2–13.9) & 9.9 (7.1–13.6) & 9.9 (7.1–13.8) & 9.9 (7.1–13.6) \\ 
Menabe & 286 / 291 & 6.7 (3.1–13.7) & 7.1 (3.1–15.2) & 6.6 (3.1–13.6) & 6.7 (3.1–13.6) & 6.7 (3.2–13.6) \\ 
Diana & 205 / 210 & 7.8 (4.8–12.5) & 7.9 (4.8–12.6) & 7.8 (4.8–12.4) & 7.9 (4.8–12.5) & 7.8 (4.8–12.4) \\ 
Sava & 219 / 232 & 9.5 (6.0–14.8) & 9.7 (6.1–15.1) & 9.4 (6.0–14.6) & 9.5 (6.0–14.7) & 9.5 (6.0–14.6) \\ 
\midrule\addlinespace[2.5pt]
\multicolumn{7}{l}{Wasting — Wealth quintile} \\[2.5pt] 
\midrule\addlinespace[2.5pt]
Poorest & 1783 / 1832 & 9.9 (8.2–11.9) & 10.1 (8.3–12.2) & 9.8 (8.1–11.8) & 9.9 (8.2–11.9) & 9.9 (8.2–11.9) \\ 
Poorer & 1447 / 1495 & 8.0 (6.5–9.8) & 8.2 (6.6–10.0) & 7.9 (6.4–9.7) & 8.1 (6.6–9.9) & 8.0 (6.5–9.9) \\ 
Middle & 1292 / 1324 & 6.9 (5.3–8.9) & 6.7 (5.1–8.6) & 7.1 (5.4–9.1) & 7.0 (5.4–8.9) & 6.8 (5.3–8.8) \\ 
Richer & 993 / 1024 & 5.5 (4.1–7.3) & 5.5 (4.1–7.4) & 5.5 (4.1–7.3) & 5.5 (4.1–7.3) & 5.5 (4.1–7.3) \\ 
Richest & 963 / 1007 & 7.1 (5.2–9.5) & 6.9 (5.1–9.3) & 7.2 (5.3–9.7) & 7.0 (5.2–9.5) & 7.0 (5.2–9.4) \\ 
\bottomrule
\end{longtable}
\begin{minipage}{\linewidth}
\vspace{.05em}
CCA, complete-case analysis (DHS definition); CCA (outcome + covariates), additionally requires observed mother's age and education; IPW, inverse probability weighting; MI-PMM, multiple imputation by predictive mean matching; MI-RF, random-forest multiple imputation (m = 20). All estimates survey-weighted; CIs on the logit scale. $\dagger$ Top-5 worst-affected status differs between methods.\\
\end{minipage}
\endgroup

\end{landscape}
\clearpage

% --- table3_simulation
\begingroup
\fontsize{8.0pt}{10.0pt}\selectfont
\setlength{\LTpost}{0mm}
\begin{longtable}{l|rrrrrr}
\caption*{
{\fontsize{20}{25}\selectfont  Table 3. Simulation study: bias, RMSE and 95\% CI coverage of national prevalence estimates\fontsize{8}{10}\selectfont }
} \\ 
\toprule
 & \multicolumn{3}{c}{{Stunting (true 39.64\%)}} & \multicolumn{3}{c}{{Wasting (true 7.62\%)}} \\ 
\cmidrule(lr){2-4} \cmidrule(lr){5-7}
 & Bias (pp) & RMSE (pp) & Coverage (\%) & Bias (pp) & RMSE (pp) & Coverage (\%) \\ 
\midrule\endhead\addlinespace[2.5pt]
\multicolumn{7}{l}{MCAR, 10\% missing} \\[2.5pt] 
\midrule\addlinespace[2.5pt]
Complete data & -0.04 & 1.01 & 92.8 & -0.02 & 0.44 & 95.2 \\ 
CCA & -0.05 & 1.04 & 92.8 & -0.02 & 0.45 & 93.8 \\ 
IPW & -0.05 & 1.04 & 93.2 & -0.02 & 0.45 & 94.4 \\ 
MI-PMM & -0.03 & 1.02 & 92.4 & -0.02 & 0.44 & 94.8 \\ 
MI-RF & -0.06 & 1.04 & 92.2 & -0.08 & 0.46 & 93.8 \\ 
\midrule\addlinespace[2.5pt]
\multicolumn{7}{l}{MCAR, 30\% missing} \\[2.5pt] 
\midrule\addlinespace[2.5pt]
Complete data & -0.03 & 0.93 & 94.4 & -0.03 & 0.47 & 92.8 \\ 
CCA & 0.01 & 1.06 & 95.6 & -0.05 & 0.57 & 92.0 \\ 
IPW & 0.00 & 1.05 & 95.4 & -0.05 & 0.57 & 92.2 \\ 
MI-PMM & 0.07 & 1.02 & 93.6 & -0.03 & 0.54 & 92.8 \\ 
MI-RF & -0.05 & 1.04 & 94.6 & -0.22 & 0.59 & 89.2 \\ 
\midrule\addlinespace[2.5pt]
\multicolumn{7}{l}{MCAR, 50\% missing} \\[2.5pt] 
\midrule\addlinespace[2.5pt]
Complete data & -0.05 & 0.88 & 96.2 & 0.01 & 0.47 & 93.6 \\ 
CCA & -0.07 & 1.12 & 96.6 & 0.00 & 0.61 & 95.2 \\ 
IPW & -0.04 & 1.10 & 97.0 & -0.01 & 0.61 & 96.0 \\ 
MI-PMM & 0.09 & 1.07 & 94.4 & 0.05 & 0.58 & 92.0 \\ 
MI-RF & -0.14 & 1.08 & 93.8 & -0.27 & 0.63 & 90.6 \\ 
\midrule\addlinespace[2.5pt]
\multicolumn{7}{l}{MAR, 10\% missing} \\[2.5pt] 
\midrule\addlinespace[2.5pt]
Complete data & -0.03 & 1.00 & 94.0 & 0.05 & 0.43 & 95.8 \\ 
CCA & -0.01 & 1.02 & 94.4 & 0.01 & 0.44 & 96.0 \\ 
IPW & -0.02 & 1.02 & 94.0 & 0.04 & 0.44 & 95.4 \\ 
MI-PMM & -0.06 & 1.01 & 94.2 & 0.07 & 0.44 & 95.2 \\ 
MI-RF & -0.06 & 1.02 & 93.8 & 0.00 & 0.44 & 96.0 \\ 
\midrule\addlinespace[2.5pt]
\multicolumn{7}{l}{MAR, 30\% missing} \\[2.5pt] 
\midrule\addlinespace[2.5pt]
Complete data & 0.01 & 0.92 & 95.2 & 0.00 & 0.43 & 95.4 \\ 
CCA & 0.07 & 1.05 & 94.6 & -0.07 & 0.49 & 95.6 \\ 
IPW & 0.01 & 1.07 & 94.8 & 0.00 & 0.51 & 95.2 \\ 
MI-PMM & -0.04 & 1.01 & 95.2 & 0.09 & 0.50 & 94.4 \\ 
MI-RF & -0.09 & 1.04 & 93.4 & -0.15 & 0.51 & 94.0 \\ 
\midrule\addlinespace[2.5pt]
\multicolumn{7}{l}{MAR, 50\% missing} \\[2.5pt] 
\midrule\addlinespace[2.5pt]
Complete data & 0.02 & 0.94 & 94.8 & -0.01 & 0.43 & 96.6 \\ 
CCA & 0.09 & 1.18 & 95.4 & -0.09 & 0.59 & 95.6 \\ 
IPW & 0.02 & 1.32 & 94.4 & -0.02 & 0.64 & 95.2 \\ 
MI-PMM & 0.00 & 1.21 & 92.2 & 0.14 & 0.59 & 94.6 \\ 
MI-RF & -0.10 & 1.20 & 88.6 & -0.26 & 0.61 & 91.2 \\ 
\midrule\addlinespace[2.5pt]
\multicolumn{7}{l}{MAR-strong, 10\% missing} \\[2.5pt] 
\midrule\addlinespace[2.5pt]
Complete data & -0.01 & 0.92 & 94.6 & -0.02 & 0.45 & 95.8 \\ 
CCA & -0.73 & 1.20 & 86.4 & -0.30 & 0.55 & 89.2 \\ 
IPW & 0.00 & 0.96 & 95.6 & -0.02 & 0.47 & 96.0 \\ 
MI-PMM & -0.16 & 0.98 & 95.0 & -0.10 & 0.47 & 94.4 \\ 
MI-RF & -0.24 & 0.99 & 94.0 & -0.19 & 0.50 & 91.8 \\ 
\midrule\addlinespace[2.5pt]
\multicolumn{7}{l}{MAR-strong, 30\% missing} \\[2.5pt] 
\midrule\addlinespace[2.5pt]
Complete data & 0.00 & 0.92 & 94.8 & 0.02 & 0.46 & 94.0 \\ 
CCA & -2.18 & 2.41 & 44.2 & -0.70 & 0.88 & 71.2 \\ 
IPW & -0.01 & 1.07 & 96.0 & 0.01 & 0.58 & 94.2 \\ 
MI-PMM & -0.52 & 1.17 & 91.8 & -0.20 & 0.59 & 90.2 \\ 
MI-RF & -0.86 & 1.35 & 84.2 & -0.49 & 0.72 & 79.4 \\ 
\midrule\addlinespace[2.5pt]
\multicolumn{7}{l}{MAR-strong, 50\% missing} \\[2.5pt] 
\midrule\addlinespace[2.5pt]
Complete data & 0.01 & 0.88 & 95.4 & 0.00 & 0.44 & 96.2 \\ 
CCA & -3.68 & 3.87 & 11.0 & -1.11 & 1.26 & 54.2 \\ 
IPW & -0.04 & 1.26 & 94.8 & -0.04 & 0.71 & 95.2 \\ 
MI-PMM & -0.90 & 1.52 & 86.4 & -0.33 & 0.74 & 86.0 \\ 
MI-RF & -1.69 & 2.05 & 63.4 & -0.88 & 1.06 & 59.8 \\ 
\midrule\addlinespace[2.5pt]
\multicolumn{7}{l}{MNAR, 10\% missing} \\[2.5pt] 
\midrule\addlinespace[2.5pt]
Complete data & -0.03 & 0.96 & 94.4 & 0.00 & 0.44 & 95.4 \\ 
CCA & -2.31 & 2.51 & 33.6 & -1.19 & 1.26 & 25.4 \\ 
IPW & -2.12 & 2.34 & 41.6 & -1.18 & 1.26 & 25.2 \\ 
MI-PMM & -2.16 & 2.36 & 38.8 & -1.12 & 1.20 & 28.0 \\ 
MI-RF & -1.99 & 2.21 & 45.2 & -1.14 & 1.22 & 29.2 \\ 
\midrule\addlinespace[2.5pt]
\multicolumn{7}{l}{MNAR, 30\% missing} \\[2.5pt] 
\midrule\addlinespace[2.5pt]
Complete data & 0.07 & 0.94 & 94.2 & 0.02 & 0.45 & 95.2 \\ 
CCA & -6.25 & 6.33 & 0.0 & -2.57 & 2.61 & 0.0 \\ 
IPW & -5.94 & 6.03 & 0.0 & -2.65 & 2.69 & 0.0 \\ 
MI-PMM & -6.04 & 6.12 & 0.0 & -2.47 & 2.51 & 0.0 \\ 
MI-RF & -5.69 & 5.78 & 0.0 & -2.60 & 2.63 & 0.0 \\ 
\midrule\addlinespace[2.5pt]
\multicolumn{7}{l}{MNAR, 50\% missing} \\[2.5pt] 
\midrule\addlinespace[2.5pt]
Complete data & -0.06 & 0.97 & 94.2 & 0.01 & 0.45 & 95.6 \\ 
CCA & -10.28 & 10.35 & 0.0 & -3.52 & 3.55 & 0.0 \\ 
IPW & -9.81 & 9.89 & 0.0 & -3.65 & 3.69 & 0.0 \\ 
MI-PMM & -9.90 & 9.97 & 0.0 & -3.41 & 3.44 & 0.0 \\ 
MI-RF & -9.61 & 9.68 & 0.0 & -3.64 & 3.67 & 0.0 \\ 
\bottomrule
\end{longtable}
\begin{minipage}{\linewidth}
\vspace{.05em}
500 replicates per scenario (Rao–Wu PSU bootstrap of a pseudo-population of 6,247 children, then imposed missingness). pp, percentage points. Complete data: estimate before deletion (benchmark). Maximum Monte Carlo SE across missing-data methods: bias 0.06 pp, coverage 2.2 percentage points. MI m = 10.\\
\end{minipage}
\endgroup

\clearpage

% --- tableS2_prevalence_sensitivity
\begin{landscape}
\begingroup
\fontsize{8.0pt}{10.0pt}\selectfont
\setlength{\LTpost}{0mm}
\begin{longtable}{l|rlllll}
\caption*{
{\fontsize{20}{25}\selectfont  Table S2. Prevalence by missing-data method, sensitivity analysis excluding WHO-flagged z-scores\fontsize{8}{10}\selectfont }
} \\ 
\toprule
 &  & \multicolumn{5}{c}{{Prevalence, \% (95\% CI)}} \\ 
\cmidrule(lr){3-7}
 & n valid / n & CCA & CCA (outcome + covariates) & IPW & MI-PMM & MI-RF \\ 
\midrule\endhead\addlinespace[2.5pt]
\multicolumn{7}{l}{Stunting — National} \\[2.5pt] 
\midrule\addlinespace[2.5pt]
National & 6402 / 6659 & 39.8 (38.0–41.7) & 39.7 (37.9–41.6) & 39.8 (38.0–41.7) & 39.8 (38.1–41.6) & 39.9 (38.1–41.7) \\ 
\midrule\addlinespace[2.5pt]
\multicolumn{7}{l}{Stunting — Region} \\[2.5pt] 
\midrule\addlinespace[2.5pt]
Antananarivo & 300 / 318 & 44.0 (39.1–49.1) & 43.4 (38.3–48.7) & 44.5 (39.6–49.5) & 44.1 (39.0–49.2) & 43.5 (38.4–48.7) \\ 
Analamanga (excl. Antananarivo) $\dagger$ & 217 / 221 & 48.3 (40.2–56.5) & 48.3 (40.2–56.4) & 48.2 (40.1–56.5) & 48.2 (40.0–56.5) & 48.3 (40.2–56.4) \\ 
Vakinankaratra & 263 / 274 & 51.9 (43.3–60.4) & 52.1 (43.3–60.8) & 51.5 (43.0–59.9) & 51.8 (43.3–60.2) & 51.6 (43.0–60.0) \\ 
Itasy & 258 / 261 & 51.6 (42.7–60.3) & 51.5 (42.8–60.2) & 51.6 (42.7–60.3) & 51.6 (42.8–60.3) & 51.4 (42.6–60.2) \\ 
Bongolava & 314 / 325 & 37.6 (30.6–45.2) & 37.5 (30.4–45.3) & 37.6 (30.6–45.2) & 37.7 (30.7–45.2) & 37.6 (30.7–45.1) \\ 
Haute Matsiatra & 250 / 259 & 49.2 (41.9–56.5) & 49.2 (41.9–56.5) & 49.0 (41.7–56.4) & 49.4 (42.2–56.7) & 49.2 (41.9–56.5) \\ 
Amoron'i Mania & 280 / 286 & 49.1 (41.6–56.6) & 49.0 (41.4–56.7) & 49.6 (42.0–57.2) & 49.1 (41.7–56.7) & 49.0 (41.6–56.5) \\ 
Vatovavy Fitovinany & 261 / 289 & 45.7 (38.5–53.1) & 45.8 (38.4–53.3) & 46.1 (38.7–53.7) & 45.5 (38.5–52.6) & 45.6 (38.5–52.8) \\ 
Ihorombe & 318 / 321 & 32.6 (26.0–40.0) & 32.7 (26.1–40.1) & 32.6 (26.1–40.0) & 32.6 (26.0–39.9) & 32.7 (26.2–40.0) \\ 
Atsimo Atsinanana $\dagger$ & 366 / 375 & 48.5 (43.1–54.0) & 48.3 (42.8–53.8) & 48.5 (43.1–54.0) & 48.5 (43.1–53.9) & 48.5 (43.1–54.0) \\ 
Atsinanana & 259 / 270 & 37.6 (31.1–44.6) & 37.0 (30.5–43.9) & 38.5 (31.4–46.1) & 37.8 (31.4–44.7) & 38.0 (31.5–44.9) \\ 
Analanjirofo & 287 / 294 & 36.2 (28.5–44.6) & 36.2 (28.4–44.8) & 36.4 (28.7–44.8) & 36.1 (28.6–44.4) & 36.2 (28.5–44.6) \\ 
Alaotra Mangoro & 250 / 258 & 44.9 (37.8–52.1) & 45.0 (37.7–52.5) & 44.8 (37.9–51.9) & 45.1 (38.1–52.3) & 44.9 (38.0–52.0) \\ 
Boeny & 217 / 221 & 28.8 (22.1–36.6) & 28.4 (21.8–36.1) & 29.0 (22.2–36.8) & 28.7 (22.0–36.5) & 28.9 (22.3–36.6) \\ 
Sofia & 218 / 224 & 28.1 (20.8–36.9) & 27.3 (20.0–36.1) & 28.4 (21.0–37.1) & 28.2 (21.0–36.7) & 28.3 (21.0–36.9) \\ 
Betsiboka & 261 / 270 & 29.6 (23.8–36.1) & 29.9 (24.1–36.3) & 29.4 (23.4–36.2) & 30.0 (24.2–36.6) & 29.9 (24.0–36.6) \\ 
Melaky & 267 / 275 & 22.3 (19.0–25.9) & 22.0 (18.9–25.6) & 22.4 (19.1–26.2) & 22.4 (19.0–26.3) & 22.5 (19.1–26.1) \\ 
Atsimo Andrefana & 363 / 383 & 26.7 (20.7–33.8) & 26.5 (20.7–33.3) & 27.1 (20.8–34.5) & 26.7 (20.8–33.6) & 26.8 (20.8–33.9) \\ 
Androy & 429 / 473 & 45.3 (40.6–50.1) & 45.7 (40.8–50.8) & 43.3 (38.0–48.7) & 45.2 (40.3–50.1) & 45.6 (40.8–50.6) \\ 
Anosy & 327 / 332 & 46.8 (38.5–55.3) & 46.7 (38.5–55.1) & 46.9 (38.6–55.4) & 46.8 (38.5–55.2) & 47.0 (38.7–55.4) \\ 
Menabe & 284 / 289 & 34.9 (29.6–40.6) & 34.1 (28.6–40.1) & 34.9 (29.7–40.6) & 34.9 (29.6–40.6) & 34.9 (29.7–40.5) \\ 
Diana & 203 / 210 & 21.7 (15.0–30.3) & 21.2 (14.5–29.9) & 21.7 (15.0–30.4) & 21.8 (15.2–30.4) & 21.9 (15.1–30.5) \\ 
Sava & 210 / 231 & 31.0 (24.6–38.1) & 31.2 (24.8–38.3) & 31.1 (24.6–38.5) & 31.1 (24.6–38.6) & 31.2 (24.5–38.8) \\ 
\midrule\addlinespace[2.5pt]
\multicolumn{7}{l}{Stunting — Wealth quintile} \\[2.5pt] 
\midrule\addlinespace[2.5pt]
Poorest & 1749 / 1827 & 42.8 (39.8–45.9) & 42.8 (39.8–46.0) & 42.2 (39.2–45.3) & 42.8 (39.7–45.8) & 42.7 (39.7–45.8) \\ 
Poorer & 1424 / 1493 & 43.4 (40.3–46.6) & 43.2 (40.1–46.4) & 43.8 (40.6–47.1) & 43.3 (40.1–46.5) & 43.4 (40.3–46.6) \\ 
Middle & 1285 / 1316 & 40.1 (36.6–43.8) & 39.8 (36.2–43.4) & 40.3 (36.7–44.0) & 40.1 (36.5–43.7) & 40.1 (36.5–43.8) \\ 
Richer & 986 / 1021 & 39.4 (35.5–43.4) & 39.7 (35.8–43.7) & 39.2 (35.3–43.3) & 39.4 (35.5–43.5) & 39.4 (35.5–43.3) \\ 
Richest & 958 / 1002 & 30.0 (25.7–34.7) & 29.6 (25.3–34.3) & 30.0 (25.8–34.7) & 30.2 (25.9–34.8) & 30.1 (25.8–34.7) \\ 
\midrule\addlinespace[2.5pt]
\multicolumn{7}{l}{Wasting — National} \\[2.5pt] 
\midrule\addlinespace[2.5pt]
National & 6478 / 6656 & 7.7 (6.9–8.6) & 7.7 (6.9–8.6) & 7.7 (6.9–8.6) & 7.7 (6.9–8.6) & 7.7 (6.9–8.6) \\ 
\midrule\addlinespace[2.5pt]
\multicolumn{7}{l}{Wasting — Region} \\[2.5pt] 
\midrule\addlinespace[2.5pt]
Antananarivo & 300 / 318 & 3.7 (2.1–6.5) & 3.8 (2.1–6.6) & 3.5 (2.0–6.3) & 3.8 (2.1–6.6) & 3.8 (2.1–6.6) \\ 
Analamanga (excl. Antananarivo) & 215 / 219 & 5.1 (2.6–9.7) & 4.7 (2.3–9.4) & 5.3 (2.7–10.1) & 5.0 (2.6–9.6) & 5.1 (2.6–9.7) \\ 
Vakinankaratra & 261 / 272 & 4.1 (2.0–8.4) & 4.2 (2.0–8.6) & 4.1 (2.0–8.4) & 4.1 (2.0–8.3) & 4.2 (2.0–8.4) \\ 
Itasy & 257 / 260 & 3.9 (2.2–6.8) & 4.0 (2.3–6.9) & 4.0 (2.3–6.9) & 4.0 (2.3–6.9) & 4.1 (2.3–7.1) \\ 
Bongolava & 314 / 323 & 7.4 (5.0–10.7) & 7.3 (4.8–10.9) & 7.5 (5.2–10.9) & 7.4 (5.1–10.8) & 7.4 (5.0–10.7) \\ 
Haute Matsiatra & 256 / 258 & 2.4 (0.6–9.2) & 2.5 (0.6–9.4) & 2.4 (0.6–9.1) & 2.4 (0.6–9.1) & 2.4 (0.6–9.1) \\ 
Amoron'i Mania & 279 / 285 & 2.6 (1.4–4.9) & 2.7 (1.4–5.0) & 2.6 (1.4–4.8) & 2.7 (1.4–5.0) & 2.6 (1.4–4.9) \\ 
Vatovavy Fitovinany $\dagger$ & 271 / 288 & 10.2 (7.4–13.9) & 9.5 (6.5–13.7) & 10.3 (7.4–14.2) & 10.2 (7.3–14.0) & 10.0 (7.2–13.7) \\ 
Ihorombe & 318 / 321 & 5.2 (2.9–9.1) & 5.2 (2.9–9.2) & 5.2 (2.9–9.1) & 5.1 (2.9–9.1) & 5.2 (2.9–9.1) \\ 
Atsimo Atsinanana & 368 / 375 & 12.4 (9.5–15.9) & 12.5 (9.6–16.1) & 12.3 (9.5–15.8) & 12.3 (9.5–15.8) & 12.3 (9.5–15.8) \\ 
Atsinanana & 264 / 271 & 8.6 (5.1–14.2) & 8.6 (5.0–14.3) & 8.6 (5.1–14.0) & 8.6 (5.1–14.1) & 8.6 (5.1–14.1) \\ 
Analanjirofo & 288 / 293 & 11.3 (7.3–17.0) & 11.1 (7.2–16.8) & 11.3 (7.3–17.1) & 11.2 (7.3–16.9) & 11.1 (7.2–16.8) \\ 
Alaotra Mangoro & 251 / 258 & 3.7 (1.9–7.2) & 3.8 (2.0–7.3) & 3.7 (1.9–7.1) & 3.8 (1.9–7.2) & 3.8 (1.9–7.2) \\ 
Boeny & 217 / 221 & 12.4 (8.3–18.2) & 12.7 (8.5–18.4) & 12.4 (8.3–18.1) & 12.3 (8.2–18.0) & 12.4 (8.3–18.1) \\ 
Sofia & 218 / 224 & 9.8 (6.1–15.2) & 10.0 (6.3–15.4) & 9.6 (6.0–15.0) & 9.7 (6.1–15.1) & 9.8 (6.2–15.1) \\ 
Betsiboka & 262 / 271 & 7.3 (4.2–12.5) & 6.9 (4.2–11.2) & 7.7 (4.2–13.7) & 7.4 (4.2–12.7) & 7.4 (4.3–12.5) \\ 
Melaky & 270 / 275 & 5.9 (3.1–10.9) & 6.0 (3.2–11.1) & 5.8 (3.0–10.9) & 5.9 (3.1–10.9) & 5.9 (3.1–10.9) \\ 
Atsimo Andrefana & 370 / 387 & 6.8 (4.2–11.0) & 6.7 (4.0–11.0) & 6.9 (4.3–11.0) & 6.9 (4.3–11.0) & 6.7 (4.1–10.8) \\ 
Androy & 461 / 474 & 15.2 (10.8–20.8) & 15.6 (11.2–21.3) & 14.9 (10.6–20.4) & 15.1 (10.8–20.7) & 15.1 (10.8–20.8) \\ 
Anosy $\dagger$ & 328 / 332 & 9.9 (7.1–13.7) & 10.1 (7.2–13.9) & 9.9 (7.1–13.6) & 9.8 (7.1–13.5) & 9.8 (7.1–13.6) \\ 
Menabe & 286 / 290 & 6.7 (3.1–13.7) & 7.1 (3.1–15.2) & 6.6 (3.1–13.6) & 6.7 (3.1–13.7) & 6.8 (3.2–13.7) \\ 
Diana & 205 / 209 & 7.8 (4.8–12.5) & 7.9 (4.8–12.6) & 7.8 (4.8–12.5) & 7.8 (4.8–12.5) & 7.8 (4.7–12.4) \\ 
Sava & 219 / 232 & 9.5 (6.0–14.8) & 9.7 (6.1–15.1) & 9.5 (6.0–14.6) & 9.6 (6.1–14.8) & 9.5 (6.0–14.6) \\ 
\midrule\addlinespace[2.5pt]
\multicolumn{7}{l}{Wasting — Wealth quintile} \\[2.5pt] 
\midrule\addlinespace[2.5pt]
Poorest & 1783 / 1829 & 9.9 (8.2–11.9) & 10.1 (8.3–12.2) & 9.8 (8.1–11.8) & 9.9 (8.2–11.9) & 9.9 (8.2–11.9) \\ 
Poorer & 1447 / 1490 & 8.0 (6.5–9.8) & 8.2 (6.6–10.0) & 7.9 (6.4–9.7) & 8.0 (6.5–9.8) & 7.9 (6.5–9.8) \\ 
Middle & 1292 / 1312 & 6.9 (5.3–8.9) & 6.7 (5.1–8.6) & 7.1 (5.4–9.1) & 6.9 (5.3–8.9) & 6.9 (5.3–8.8) \\ 
Richer & 993 / 1021 & 5.5 (4.1–7.3) & 5.5 (4.1–7.4) & 5.5 (4.1–7.3) & 5.5 (4.1–7.4) & 5.5 (4.1–7.3) \\ 
Richest & 963 / 1004 & 7.1 (5.2–9.5) & 6.9 (5.1–9.3) & 7.1 (5.2–9.6) & 7.0 (5.2–9.5) & 7.1 (5.2–9.5) \\ 
\bottomrule
\end{longtable}
\begin{minipage}{\linewidth}
\vspace{.05em}
CCA, complete-case analysis (DHS definition); CCA (outcome + covariates), additionally requires observed mother's age and education; IPW, inverse probability weighting; MI-PMM, multiple imputation by predictive mean matching; MI-RF, random-forest multiple imputation (m = 20). All estimates survey-weighted; CIs on the logit scale. $\dagger$ Top-5 worst-affected status differs between methods.\\
\end{minipage}
\endgroup

\end{landscape}
\clearpage

% --- tableS3_simulation_full
\begin{landscape}
\begingroup
\fontsize{8.0pt}{10.0pt}\selectfont
\begin{longtable}{l|rrrrrrrrrl}
\caption*{
{\fontsize{20}{25}\selectfont  Table S3. Simulation study: full performance measures (percentage points unless stated)\fontsize{8}{10}\selectfont }
} \\ 
\toprule
 & R & Bias (MCSE) & Rel. bias (\%) & Empirical SE (MCSE) & Model SE (MCSE) & SE ratio & RMSE (MCSE) & Coverage \% (MCSE) & CI width & FMI \\ 
\midrule\endhead\addlinespace[2.5pt]
\multicolumn{11}{l}{Stunting — MCAR 10\%} \\[2.5pt] 
\midrule\addlinespace[2.5pt]
Complete data & 500 & -0.04 (0.04) & -0.1 & 1.01 (0.03) & 0.93 (0.00) & 0.92 & 1.01 (0.03) & 92.8 (1.2) & 3.63 & — \\ 
CCA & 500 & -0.05 (0.05) & -0.1 & 1.04 (0.03) & 0.96 (0.00) & 0.92 & 1.04 (0.03) & 92.8 (1.2) & 3.76 & — \\ 
IPW & 500 & -0.05 (0.05) & -0.1 & 1.04 (0.03) & 0.96 (0.00) & 0.92 & 1.04 (0.03) & 93.2 (1.1) & 3.76 & — \\ 
MI-PMM & 500 & -0.03 (0.05) & -0.1 & 1.02 (0.03) & 0.93 (0.00) & 0.92 & 1.02 (0.03) & 92.4 (1.2) & 3.67 & 0.07 \\ 
MI-RF & 500 & -0.06 (0.05) & -0.2 & 1.04 (0.03) & 0.94 (0.00) & 0.90 & 1.04 (0.03) & 92.2 (1.2) & 3.67 & 0.07 \\ 
\midrule\addlinespace[2.5pt]
\multicolumn{11}{l}{Stunting — MCAR 30\%} \\[2.5pt] 
\midrule\addlinespace[2.5pt]
Complete data & 500 & -0.03 (0.04) & -0.1 & 0.93 (0.03) & 0.93 (0.00) & 1.00 & 0.93 (0.03) & 94.4 (1.0) & 3.64 & — \\ 
CCA & 500 & 0.01 (0.05) & 0.0 & 1.06 (0.03) & 1.05 (0.00) & 0.99 & 1.06 (0.03) & 95.6 (0.9) & 4.11 & — \\ 
IPW & 500 & 0.00 (0.05) & 0.0 & 1.05 (0.03) & 1.05 (0.00) & 0.99 & 1.05 (0.03) & 95.4 (0.9) & 4.11 & — \\ 
MI-PMM & 500 & 0.07 (0.05) & 0.2 & 1.02 (0.03) & 0.98 (0.00) & 0.96 & 1.02 (0.03) & 93.6 (1.1) & 3.86 & 0.23 \\ 
MI-RF & 500 & -0.05 (0.05) & -0.1 & 1.04 (0.03) & 0.97 (0.00) & 0.93 & 1.04 (0.03) & 94.6 (1.0) & 3.81 & 0.20 \\ 
\midrule\addlinespace[2.5pt]
\multicolumn{11}{l}{Stunting — MCAR 50\%} \\[2.5pt] 
\midrule\addlinespace[2.5pt]
Complete data & 500 & -0.05 (0.04) & -0.1 & 0.88 (0.03) & 0.93 (0.00) & 1.05 & 0.88 (0.03) & 96.2 (0.9) & 3.63 & — \\ 
CCA & 500 & -0.07 (0.05) & -0.2 & 1.12 (0.04) & 1.18 (0.00) & 1.06 & 1.12 (0.03) & 96.6 (0.8) & 4.63 & — \\ 
IPW & 500 & -0.04 (0.05) & -0.1 & 1.10 (0.03) & 1.18 (0.00) & 1.07 & 1.10 (0.03) & 97.0 (0.8) & 4.64 & — \\ 
MI-PMM & 500 & 0.09 (0.05) & 0.2 & 1.07 (0.03) & 1.07 (0.00) & 1.00 & 1.07 (0.03) & 94.4 (1.0) & 4.28 & 0.41 \\ 
MI-RF & 500 & -0.14 (0.05) & -0.4 & 1.07 (0.03) & 1.00 (0.00) & 0.93 & 1.08 (0.03) & 93.8 (1.1) & 3.98 & 0.32 \\ 
\midrule\addlinespace[2.5pt]
\multicolumn{11}{l}{Stunting — MAR 10\%} \\[2.5pt] 
\midrule\addlinespace[2.5pt]
Complete data & 500 & -0.03 (0.04) & -0.1 & 1.00 (0.03) & 0.92 (0.00) & 0.92 & 1.00 (0.04) & 94.0 (1.1) & 3.63 & — \\ 
CCA & 500 & -0.01 (0.05) & 0.0 & 1.02 (0.03) & 0.96 (0.00) & 0.94 & 1.02 (0.04) & 94.4 (1.0) & 3.77 & — \\ 
IPW & 500 & -0.02 (0.05) & -0.1 & 1.02 (0.03) & 0.96 (0.00) & 0.94 & 1.02 (0.04) & 94.0 (1.1) & 3.77 & — \\ 
MI-PMM & 500 & -0.06 (0.05) & -0.2 & 1.01 (0.03) & 0.94 (0.00) & 0.93 & 1.01 (0.04) & 94.2 (1.0) & 3.68 & 0.08 \\ 
MI-RF & 500 & -0.06 (0.05) & -0.1 & 1.02 (0.03) & 0.94 (0.00) & 0.91 & 1.02 (0.04) & 93.8 (1.1) & 3.67 & 0.07 \\ 
\midrule\addlinespace[2.5pt]
\multicolumn{11}{l}{Stunting — MAR 30\%} \\[2.5pt] 
\midrule\addlinespace[2.5pt]
Complete data & 500 & 0.01 (0.04) & 0.0 & 0.92 (0.03) & 0.93 (0.00) & 1.00 & 0.92 (0.03) & 95.2 (1.0) & 3.64 & — \\ 
CCA & 500 & 0.07 (0.05) & 0.2 & 1.05 (0.03) & 1.05 (0.00) & 1.00 & 1.05 (0.03) & 94.6 (1.0) & 4.13 & — \\ 
IPW & 500 & 0.01 (0.05) & 0.0 & 1.07 (0.03) & 1.08 (0.00) & 1.01 & 1.07 (0.03) & 94.8 (1.0) & 4.25 & — \\ 
MI-PMM & 500 & -0.04 (0.05) & -0.1 & 1.01 (0.03) & 1.00 (0.00) & 0.99 & 1.01 (0.03) & 95.2 (1.0) & 3.95 & 0.26 \\ 
MI-RF & 500 & -0.09 (0.05) & -0.2 & 1.04 (0.03) & 0.97 (0.00) & 0.93 & 1.04 (0.03) & 93.4 (1.1) & 3.83 & 0.21 \\ 
\midrule\addlinespace[2.5pt]
\multicolumn{11}{l}{Stunting — MAR 30\% (ntree=50)} \\[2.5pt] 
\midrule\addlinespace[2.5pt]
MI-RF (ntree=50) & 200 & -0.10 (0.07) & -0.3 & 0.97 (0.05) & 0.97 (0.00) & 0.99 & 0.97 (0.05) & 96.5 (1.3) & 3.81 & 0.20 \\ 
\midrule\addlinespace[2.5pt]
\multicolumn{11}{l}{Stunting — MAR 50\%} \\[2.5pt] 
\midrule\addlinespace[2.5pt]
Complete data & 500 & 0.02 (0.04) & 0.0 & 0.94 (0.03) & 0.93 (0.00) & 0.98 & 0.94 (0.03) & 94.8 (1.0) & 3.64 & — \\ 
CCA & 500 & 0.09 (0.05) & 0.2 & 1.18 (0.04) & 1.19 (0.00) & 1.01 & 1.18 (0.04) & 95.4 (0.9) & 4.68 & — \\ 
IPW & 500 & 0.02 (0.06) & 0.0 & 1.32 (0.04) & 1.30 (0.00) & 0.98 & 1.32 (0.04) & 94.4 (1.0) & 5.09 & — \\ 
MI-PMM & 500 & 0.00 (0.05) & 0.0 & 1.21 (0.04) & 1.14 (0.01) & 0.94 & 1.21 (0.04) & 92.2 (1.2) & 4.58 & 0.47 \\ 
MI-RF & 500 & -0.10 (0.05) & -0.3 & 1.20 (0.04) & 1.01 (0.00) & 0.85 & 1.20 (0.04) & 88.6 (1.4) & 4.04 & 0.34 \\ 
\midrule\addlinespace[2.5pt]
\multicolumn{11}{l}{Stunting — MAR-strong 10\%} \\[2.5pt] 
\midrule\addlinespace[2.5pt]
Complete data & 500 & -0.01 (0.04) & 0.0 & 0.92 (0.03) & 0.92 (0.00) & 1.00 & 0.92 (0.03) & 94.6 (1.0) & 3.63 & — \\ 
CCA & 500 & -0.73 (0.04) & -1.8 & 0.96 (0.03) & 0.96 (0.00) & 0.99 & 1.20 (0.04) & 86.4 (1.5) & 3.75 & — \\ 
IPW & 500 & 0.00 (0.04) & 0.0 & 0.96 (0.03) & 0.96 (0.00) & 1.00 & 0.96 (0.03) & 95.6 (0.9) & 3.76 & — \\ 
MI-PMM & 500 & -0.16 (0.04) & -0.4 & 0.96 (0.03) & 0.93 (0.00) & 0.97 & 0.98 (0.03) & 95.0 (1.0) & 3.67 & 0.08 \\ 
MI-RF & 500 & -0.24 (0.04) & -0.6 & 0.96 (0.03) & 0.93 (0.00) & 0.97 & 0.99 (0.03) & 94.0 (1.1) & 3.66 & 0.07 \\ 
\midrule\addlinespace[2.5pt]
\multicolumn{11}{l}{Stunting — MAR-strong 30\%} \\[2.5pt] 
\midrule\addlinespace[2.5pt]
Complete data & 500 & 0.00 (0.04) & 0.0 & 0.92 (0.03) & 0.93 (0.00) & 1.01 & 0.92 (0.03) & 94.8 (1.0) & 3.64 & — \\ 
CCA & 500 & -2.18 (0.05) & -5.5 & 1.04 (0.03) & 1.04 (0.00) & 1.01 & 2.41 (0.04) & 44.2 (2.2) & 4.09 & — \\ 
IPW & 500 & -0.01 (0.05) & 0.0 & 1.07 (0.03) & 1.08 (0.00) & 1.00 & 1.07 (0.04) & 96.0 (0.9) & 4.23 & — \\ 
MI-PMM & 500 & -0.52 (0.05) & -1.3 & 1.04 (0.03) & 0.99 (0.00) & 0.95 & 1.17 (0.04) & 91.8 (1.2) & 3.92 & 0.26 \\ 
MI-RF & 500 & -0.86 (0.05) & -2.2 & 1.05 (0.03) & 0.96 (0.00) & 0.92 & 1.35 (0.04) & 84.2 (1.6) & 3.80 & 0.21 \\ 
\midrule\addlinespace[2.5pt]
\multicolumn{11}{l}{Stunting — MAR-strong 50\%} \\[2.5pt] 
\midrule\addlinespace[2.5pt]
Complete data & 500 & 0.01 (0.04) & 0.0 & 0.88 (0.03) & 0.93 (0.00) & 1.05 & 0.88 (0.03) & 95.4 (0.9) & 3.64 & — \\ 
CCA & 500 & -3.68 (0.05) & -9.3 & 1.17 (0.04) & 1.18 (0.00) & 1.01 & 3.87 (0.05) & 11.0 (1.4) & 4.62 & — \\ 
IPW & 500 & -0.04 (0.06) & -0.1 & 1.26 (0.04) & 1.30 (0.00) & 1.03 & 1.26 (0.04) & 94.8 (1.0) & 5.10 & — \\ 
MI-PMM & 500 & -0.90 (0.05) & -2.3 & 1.22 (0.04) & 1.12 (0.01) & 0.92 & 1.52 (0.05) & 86.4 (1.5) & 4.52 & 0.47 \\ 
MI-RF & 500 & -1.69 (0.05) & -4.3 & 1.16 (0.04) & 1.01 (0.00) & 0.87 & 2.05 (0.05) & 63.4 (2.2) & 4.01 & 0.35 \\ 
\midrule\addlinespace[2.5pt]
\multicolumn{11}{l}{Stunting — MNAR 10\%} \\[2.5pt] 
\midrule\addlinespace[2.5pt]
Complete data & 500 & -0.03 (0.04) & -0.1 & 0.96 (0.03) & 0.93 (0.00) & 0.97 & 0.96 (0.03) & 94.4 (1.0) & 3.65 & — \\ 
CCA & 500 & -2.31 (0.04) & -5.8 & 0.97 (0.03) & 0.95 (0.00) & 0.98 & 2.51 (0.04) & 33.6 (2.1) & 3.71 & — \\ 
IPW & 500 & -2.12 (0.04) & -5.4 & 0.97 (0.03) & 0.95 (0.00) & 0.97 & 2.34 (0.04) & 41.6 (2.2) & 3.71 & — \\ 
MI-PMM & 500 & -2.16 (0.04) & -5.4 & 0.96 (0.03) & 0.92 (0.00) & 0.96 & 2.36 (0.04) & 38.8 (2.2) & 3.63 & 0.07 \\ 
MI-RF & 500 & -1.99 (0.04) & -5.0 & 0.97 (0.03) & 0.93 (0.00) & 0.95 & 2.21 (0.04) & 45.2 (2.2) & 3.63 & 0.07 \\ 
\midrule\addlinespace[2.5pt]
\multicolumn{11}{l}{Stunting — MNAR 30\%} \\[2.5pt] 
\midrule\addlinespace[2.5pt]
Complete data & 500 & 0.07 (0.04) & 0.2 & 0.93 (0.03) & 0.92 (0.00) & 0.99 & 0.94 (0.03) & 94.2 (1.0) & 3.63 & — \\ 
CCA & 500 & -6.25 (0.05) & -15.8 & 1.02 (0.03) & 1.00 (0.00) & 0.98 & 6.33 (0.05) & 0.0 (0.0) & 3.93 & — \\ 
IPW & 500 & -5.94 (0.05) & -15.0 & 1.03 (0.03) & 1.03 (0.00) & 1.00 & 6.03 (0.05) & 0.0 (0.0) & 4.02 & — \\ 
MI-PMM & 500 & -6.04 (0.04) & -15.2 & 0.99 (0.03) & 0.95 (0.00) & 0.96 & 6.12 (0.04) & 0.0 (0.0) & 3.76 & 0.26 \\ 
MI-RF & 500 & -5.69 (0.05) & -14.4 & 1.01 (0.03) & 0.93 (0.00) & 0.92 & 5.78 (0.05) & 0.0 (0.0) & 3.68 & 0.22 \\ 
\midrule\addlinespace[2.5pt]
\multicolumn{11}{l}{Stunting — MNAR 50\%} \\[2.5pt] 
\midrule\addlinespace[2.5pt]
Complete data & 500 & -0.06 (0.04) & -0.2 & 0.96 (0.03) & 0.92 (0.00) & 0.96 & 0.97 (0.03) & 94.2 (1.0) & 3.62 & — \\ 
CCA & 500 & -10.28 (0.05) & -25.9 & 1.14 (0.04) & 1.10 (0.00) & 0.96 & 10.35 (0.05) & 0.0 (0.0) & 4.32 & — \\ 
IPW & 500 & -9.81 (0.06) & -24.8 & 1.26 (0.04) & 1.19 (0.00) & 0.95 & 9.89 (0.06) & 0.0 (0.0) & 4.67 & — \\ 
MI-PMM & 500 & -9.90 (0.05) & -25.0 & 1.17 (0.04) & 1.02 (0.01) & 0.87 & 9.97 (0.05) & 0.0 (0.0) & 4.10 & 0.45 \\ 
MI-RF & 500 & -9.61 (0.05) & -24.3 & 1.17 (0.04) & 0.94 (0.00) & 0.81 & 9.68 (0.05) & 0.0 (0.0) & 3.76 & 0.35 \\ 
\midrule\addlinespace[2.5pt]
\multicolumn{11}{l}{Wasting — MCAR 10\%} \\[2.5pt] 
\midrule\addlinespace[2.5pt]
Complete data & 500 & -0.02 (0.02) & -0.3 & 0.44 (0.01) & 0.45 (0.00) & 1.02 & 0.44 (0.01) & 95.2 (1.0) & 1.75 & — \\ 
CCA & 500 & -0.02 (0.02) & -0.3 & 0.45 (0.01) & 0.47 (0.00) & 1.02 & 0.45 (0.01) & 93.8 (1.1) & 1.83 & — \\ 
IPW & 500 & -0.02 (0.02) & -0.3 & 0.45 (0.01) & 0.47 (0.00) & 1.03 & 0.45 (0.01) & 94.4 (1.0) & 1.83 & — \\ 
MI-PMM & 500 & -0.02 (0.02) & -0.3 & 0.44 (0.01) & 0.46 (0.00) & 1.03 & 0.44 (0.01) & 94.8 (1.0) & 1.80 & 0.08 \\ 
MI-RF & 500 & -0.08 (0.02) & -1.0 & 0.45 (0.01) & 0.46 (0.00) & 1.01 & 0.46 (0.01) & 93.8 (1.1) & 1.80 & 0.08 \\ 
\midrule\addlinespace[2.5pt]
\multicolumn{11}{l}{Wasting — MCAR 30\%} \\[2.5pt] 
\midrule\addlinespace[2.5pt]
Complete data & 500 & -0.03 (0.02) & -0.4 & 0.47 (0.01) & 0.45 (0.00) & 0.95 & 0.47 (0.01) & 92.8 (1.2) & 1.76 & — \\ 
CCA & 500 & -0.05 (0.03) & -0.6 & 0.57 (0.02) & 0.52 (0.00) & 0.90 & 0.57 (0.02) & 92.0 (1.2) & 2.03 & — \\ 
IPW & 500 & -0.05 (0.03) & -0.6 & 0.57 (0.02) & 0.52 (0.00) & 0.91 & 0.57 (0.02) & 92.2 (1.2) & 2.04 & — \\ 
MI-PMM & 500 & -0.03 (0.02) & -0.4 & 0.54 (0.02) & 0.49 (0.00) & 0.90 & 0.54 (0.02) & 92.8 (1.2) & 1.93 & 0.24 \\ 
MI-RF & 500 & -0.22 (0.02) & -2.9 & 0.55 (0.02) & 0.48 (0.00) & 0.87 & 0.59 (0.02) & 89.2 (1.4) & 1.90 & 0.22 \\ 
\midrule\addlinespace[2.5pt]
\multicolumn{11}{l}{Wasting — MCAR 50\%} \\[2.5pt] 
\midrule\addlinespace[2.5pt]
Complete data & 500 & 0.01 (0.02) & 0.1 & 0.47 (0.01) & 0.45 (0.00) & 0.96 & 0.47 (0.02) & 93.6 (1.1) & 1.76 & — \\ 
CCA & 500 & 0.00 (0.03) & -0.1 & 0.61 (0.02) & 0.60 (0.00) & 0.98 & 0.61 (0.02) & 95.2 (1.0) & 2.35 & — \\ 
IPW & 500 & -0.01 (0.03) & -0.1 & 0.61 (0.02) & 0.60 (0.00) & 0.99 & 0.61 (0.02) & 96.0 (0.9) & 2.36 & — \\ 
MI-PMM & 500 & 0.05 (0.03) & 0.6 & 0.58 (0.02) & 0.52 (0.00) & 0.91 & 0.58 (0.02) & 92.0 (1.2) & 2.10 & 0.37 \\ 
MI-RF & 500 & -0.27 (0.03) & -3.5 & 0.57 (0.02) & 0.51 (0.00) & 0.90 & 0.63 (0.02) & 90.6 (1.3) & 2.04 & 0.34 \\ 
\midrule\addlinespace[2.5pt]
\multicolumn{11}{l}{Wasting — MAR 10\%} \\[2.5pt] 
\midrule\addlinespace[2.5pt]
Complete data & 500 & 0.05 (0.02) & 0.7 & 0.43 (0.01) & 0.45 (0.00) & 1.05 & 0.43 (0.01) & 95.8 (0.9) & 1.77 & — \\ 
CCA & 500 & 0.01 (0.02) & 0.2 & 0.44 (0.01) & 0.47 (0.00) & 1.06 & 0.44 (0.01) & 96.0 (0.9) & 1.83 & — \\ 
IPW & 500 & 0.04 (0.02) & 0.5 & 0.44 (0.01) & 0.47 (0.00) & 1.07 & 0.44 (0.01) & 95.4 (0.9) & 1.85 & — \\ 
MI-PMM & 500 & 0.07 (0.02) & 0.9 & 0.43 (0.01) & 0.46 (0.00) & 1.07 & 0.44 (0.01) & 95.2 (1.0) & 1.81 & 0.09 \\ 
MI-RF & 500 & 0.00 (0.02) & 0.0 & 0.44 (0.01) & 0.46 (0.00) & 1.05 & 0.44 (0.01) & 96.0 (0.9) & 1.81 & 0.09 \\ 
\midrule\addlinespace[2.5pt]
\multicolumn{11}{l}{Wasting — MAR 30\%} \\[2.5pt] 
\midrule\addlinespace[2.5pt]
Complete data & 500 & 0.00 (0.02) & 0.0 & 0.43 (0.01) & 0.45 (0.00) & 1.03 & 0.43 (0.01) & 95.4 (0.9) & 1.76 & — \\ 
CCA & 500 & -0.07 (0.02) & -0.9 & 0.49 (0.02) & 0.51 (0.00) & 1.05 & 0.49 (0.01) & 95.6 (0.9) & 2.02 & — \\ 
IPW & 500 & 0.00 (0.02) & 0.0 & 0.51 (0.02) & 0.54 (0.00) & 1.05 & 0.51 (0.02) & 95.2 (1.0) & 2.12 & — \\ 
MI-PMM & 500 & 0.09 (0.02) & 1.2 & 0.49 (0.02) & 0.50 (0.00) & 1.02 & 0.50 (0.02) & 94.4 (1.0) & 1.97 & 0.28 \\ 
MI-RF & 500 & -0.15 (0.02) & -1.9 & 0.48 (0.02) & 0.48 (0.00) & 1.00 & 0.51 (0.01) & 94.0 (1.1) & 1.92 & 0.24 \\ 
\midrule\addlinespace[2.5pt]
\multicolumn{11}{l}{Wasting — MAR 30\% (ntree=50)} \\[2.5pt] 
\midrule\addlinespace[2.5pt]
MI-RF (ntree=50) & 200 & -0.18 (0.03) & -2.4 & 0.44 (0.02) & 0.48 (0.00) & 1.09 & 0.47 (0.02) & 97.0 (1.2) & 1.88 & 0.23 \\ 
\midrule\addlinespace[2.5pt]
\multicolumn{11}{l}{Wasting — MAR 50\%} \\[2.5pt] 
\midrule\addlinespace[2.5pt]
Complete data & 500 & -0.01 (0.02) & -0.1 & 0.43 (0.01) & 0.45 (0.00) & 1.04 & 0.43 (0.01) & 96.6 (0.8) & 1.76 & — \\ 
CCA & 500 & -0.09 (0.03) & -1.2 & 0.58 (0.02) & 0.59 (0.00) & 1.02 & 0.59 (0.02) & 95.6 (0.9) & 2.32 & — \\ 
IPW & 500 & -0.02 (0.03) & -0.3 & 0.64 (0.02) & 0.66 (0.00) & 1.03 & 0.64 (0.02) & 95.2 (1.0) & 2.58 & — \\ 
MI-PMM & 500 & 0.14 (0.03) & 1.8 & 0.57 (0.02) & 0.56 (0.00) & 0.98 & 0.59 (0.02) & 94.6 (1.0) & 2.24 & 0.44 \\ 
MI-RF & 500 & -0.26 (0.02) & -3.4 & 0.55 (0.02) & 0.52 (0.00) & 0.93 & 0.61 (0.02) & 91.2 (1.3) & 2.06 & 0.37 \\ 
\midrule\addlinespace[2.5pt]
\multicolumn{11}{l}{Wasting — MAR-strong 10\%} \\[2.5pt] 
\midrule\addlinespace[2.5pt]
Complete data & 500 & -0.02 (0.02) & -0.3 & 0.45 (0.01) & 0.45 (0.00) & 1.00 & 0.45 (0.01) & 95.8 (0.9) & 1.75 & — \\ 
CCA & 500 & -0.30 (0.02) & -3.9 & 0.46 (0.01) & 0.46 (0.00) & 1.01 & 0.55 (0.02) & 89.2 (1.4) & 1.81 & — \\ 
IPW & 500 & -0.02 (0.02) & -0.2 & 0.47 (0.01) & 0.47 (0.00) & 1.00 & 0.47 (0.01) & 96.0 (0.9) & 1.85 & — \\ 
MI-PMM & 500 & -0.10 (0.02) & -1.3 & 0.46 (0.01) & 0.46 (0.00) & 0.99 & 0.47 (0.01) & 94.4 (1.0) & 1.80 & 0.10 \\ 
MI-RF & 500 & -0.19 (0.02) & -2.5 & 0.46 (0.01) & 0.46 (0.00) & 0.98 & 0.50 (0.01) & 91.8 (1.2) & 1.80 & 0.09 \\ 
\midrule\addlinespace[2.5pt]
\multicolumn{11}{l}{Wasting — MAR-strong 30\%} \\[2.5pt] 
\midrule\addlinespace[2.5pt]
Complete data & 500 & 0.02 (0.02) & 0.3 & 0.46 (0.01) & 0.45 (0.00) & 0.97 & 0.46 (0.02) & 94.0 (1.1) & 1.77 & — \\ 
CCA & 500 & -0.70 (0.02) & -9.2 & 0.53 (0.02) & 0.51 (0.00) & 0.95 & 0.88 (0.02) & 71.2 (2.0) & 1.99 & — \\ 
IPW & 500 & 0.01 (0.03) & 0.1 & 0.58 (0.02) & 0.56 (0.00) & 0.97 & 0.58 (0.02) & 94.2 (1.0) & 2.20 & — \\ 
MI-PMM & 500 & -0.20 (0.02) & -2.6 & 0.55 (0.02) & 0.49 (0.00) & 0.90 & 0.59 (0.02) & 90.2 (1.3) & 1.96 & 0.27 \\ 
MI-RF & 500 & -0.49 (0.02) & -6.5 & 0.53 (0.02) & 0.48 (0.00) & 0.90 & 0.72 (0.02) & 79.4 (1.8) & 1.90 & 0.23 \\ 
\midrule\addlinespace[2.5pt]
\multicolumn{11}{l}{Wasting — MAR-strong 50\%} \\[2.5pt] 
\midrule\addlinespace[2.5pt]
Complete data & 500 & 0.00 (0.02) & 0.0 & 0.44 (0.01) & 0.45 (0.00) & 1.03 & 0.44 (0.01) & 96.2 (0.9) & 1.76 & — \\ 
CCA & 500 & -1.11 (0.03) & -14.6 & 0.59 (0.02) & 0.57 (0.00) & 0.96 & 1.26 (0.02) & 54.2 (2.2) & 2.23 & — \\ 
IPW & 500 & -0.04 (0.03) & -0.6 & 0.71 (0.02) & 0.70 (0.00) & 1.00 & 0.71 (0.02) & 95.2 (1.0) & 2.77 & — \\ 
MI-PMM & 500 & -0.33 (0.03) & -4.3 & 0.66 (0.02) & 0.54 (0.00) & 0.81 & 0.74 (0.02) & 86.0 (1.6) & 2.16 & 0.43 \\ 
MI-RF & 500 & -0.88 (0.03) & -11.6 & 0.58 (0.02) & 0.50 (0.00) & 0.86 & 1.06 (0.02) & 59.8 (2.2) & 1.99 & 0.36 \\ 
\midrule\addlinespace[2.5pt]
\multicolumn{11}{l}{Wasting — MNAR 10\%} \\[2.5pt] 
\midrule\addlinespace[2.5pt]
Complete data & 500 & 0.00 (0.02) & 0.0 & 0.44 (0.01) & 0.45 (0.00) & 1.02 & 0.44 (0.01) & 95.4 (0.9) & 1.76 & — \\ 
CCA & 500 & -1.19 (0.02) & -15.6 & 0.43 (0.01) & 0.43 (0.00) & 0.98 & 1.26 (0.02) & 25.4 (1.9) & 1.68 & — \\ 
IPW & 500 & -1.18 (0.02) & -15.5 & 0.43 (0.01) & 0.43 (0.00) & 0.99 & 1.26 (0.02) & 25.2 (1.9) & 1.68 & — \\ 
MI-PMM & 500 & -1.12 (0.02) & -14.7 & 0.43 (0.01) & 0.42 (0.00) & 0.99 & 1.20 (0.02) & 28.0 (2.0) & 1.66 & 0.10 \\ 
MI-RF & 500 & -1.14 (0.02) & -14.9 & 0.44 (0.01) & 0.42 (0.00) & 0.97 & 1.22 (0.02) & 29.2 (2.0) & 1.67 & 0.09 \\ 
\midrule\addlinespace[2.5pt]
\multicolumn{11}{l}{Wasting — MNAR 30\%} \\[2.5pt] 
\midrule\addlinespace[2.5pt]
Complete data & 500 & 0.02 (0.02) & 0.2 & 0.45 (0.01) & 0.45 (0.00) & 1.01 & 0.45 (0.01) & 95.2 (1.0) & 1.76 & — \\ 
CCA & 500 & -2.57 (0.02) & -33.7 & 0.45 (0.01) & 0.43 (0.00) & 0.95 & 2.61 (0.02) & 0.0 (0.0) & 1.67 & — \\ 
IPW & 500 & -2.65 (0.02) & -34.8 & 0.45 (0.01) & 0.43 (0.00) & 0.96 & 2.69 (0.02) & 0.0 (0.0) & 1.69 & — \\ 
MI-PMM & 500 & -2.47 (0.02) & -32.4 & 0.43 (0.01) & 0.41 (0.00) & 0.94 & 2.51 (0.02) & 0.0 (0.0) & 1.61 & 0.26 \\ 
MI-RF & 500 & -2.60 (0.02) & -34.1 & 0.44 (0.01) & 0.40 (0.00) & 0.92 & 2.63 (0.02) & 0.0 (0.0) & 1.59 & 0.24 \\ 
\midrule\addlinespace[2.5pt]
\multicolumn{11}{l}{Wasting — MNAR 50\%} \\[2.5pt] 
\midrule\addlinespace[2.5pt]
Complete data & 500 & 0.01 (0.02) & 0.1 & 0.45 (0.01) & 0.45 (0.00) & 1.00 & 0.45 (0.01) & 95.6 (0.9) & 1.76 & — \\ 
CCA & 500 & -3.52 (0.02) & -46.2 & 0.48 (0.02) & 0.45 (0.00) & 0.93 & 3.55 (0.02) & 0.0 (0.0) & 1.77 & — \\ 
IPW & 500 & -3.65 (0.02) & -47.9 & 0.50 (0.02) & 0.47 (0.00) & 0.93 & 3.69 (0.02) & 0.0 (0.0) & 1.84 & — \\ 
MI-PMM & 500 & -3.41 (0.02) & -44.8 & 0.47 (0.01) & 0.40 (0.00) & 0.87 & 3.44 (0.02) & 0.0 (0.0) & 1.62 & 0.41 \\ 
MI-RF & 500 & -3.64 (0.02) & -47.8 & 0.46 (0.01) & 0.39 (0.00) & 0.84 & 3.67 (0.02) & 0.0 (0.0) & 1.54 & 0.35 \\ 
\bottomrule
\end{longtable}
\endgroup

\end{landscape}
\clearpage

% --- tableS4_simulation_ranking
\begingroup
\fontsize{8.0pt}{10.0pt}\selectfont
\begin{longtable}{l|rrrr}
\caption*{
{\fontsize{20}{25}\selectfont  Table S4. Simulation study: recovery of the true regional ranking (23 regions)\fontsize{8}{10}\selectfont }
} \\ 
\toprule
 & Mean Spearman $\rho$ & Mean top-5 overlap & Top-5 exactly recovered (\%) & Worst region correct (\%) \\ 
\midrule\endhead\addlinespace[2.5pt]
\multicolumn{5}{l}{Stunting — MCAR 10\%} \\[2.5pt] 
\midrule\addlinespace[2.5pt]
Complete data & 0.92 & 0.69 & 6.4 & 34.2 \\ 
CCA & 0.92 & 0.68 & 6.4 & 33.6 \\ 
IPW & 0.92 & 0.68 & 6.2 & 33.2 \\ 
MI-PMM & 0.92 & 0.68 & 4.8 & 33.0 \\ 
MI-RF & 0.92 & 0.68 & 5.6 & 34.0 \\ 
\midrule\addlinespace[2.5pt]
\multicolumn{5}{l}{Stunting — MCAR 30\%} \\[2.5pt] 
\midrule\addlinespace[2.5pt]
Complete data & 0.92 & 0.69 & 5.6 & 28.2 \\ 
CCA & 0.91 & 0.66 & 3.8 & 25.4 \\ 
IPW & 0.91 & 0.66 & 3.8 & 25.8 \\ 
MI-PMM & 0.91 & 0.66 & 3.6 & 30.4 \\ 
MI-RF & 0.91 & 0.68 & 5.2 & 31.2 \\ 
\midrule\addlinespace[2.5pt]
\multicolumn{5}{l}{Stunting — MCAR 50\%} \\[2.5pt] 
\midrule\addlinespace[2.5pt]
Complete data & 0.92 & 0.69 & 6.4 & 32.6 \\ 
CCA & 0.89 & 0.63 & 4.0 & 27.8 \\ 
IPW & 0.89 & 0.64 & 4.4 & 27.6 \\ 
MI-PMM & 0.89 & 0.64 & 3.2 & 30.4 \\ 
MI-RF & 0.89 & 0.67 & 5.2 & 33.0 \\ 
\midrule\addlinespace[2.5pt]
\multicolumn{5}{l}{Stunting — MAR 10\%} \\[2.5pt] 
\midrule\addlinespace[2.5pt]
Complete data & 0.92 & 0.70 & 7.4 & 30.0 \\ 
CCA & 0.92 & 0.71 & 8.0 & 32.0 \\ 
IPW & 0.92 & 0.69 & 6.4 & 29.4 \\ 
MI-PMM & 0.92 & 0.71 & 8.4 & 32.4 \\ 
MI-RF & 0.92 & 0.71 & 9.4 & 30.8 \\ 
\midrule\addlinespace[2.5pt]
\multicolumn{5}{l}{Stunting — MAR 30\%} \\[2.5pt] 
\midrule\addlinespace[2.5pt]
Complete data & 0.92 & 0.69 & 6.8 & 31.0 \\ 
CCA & 0.91 & 0.68 & 5.8 & 35.2 \\ 
IPW & 0.90 & 0.65 & 4.0 & 31.2 \\ 
MI-PMM & 0.91 & 0.68 & 6.0 & 35.6 \\ 
MI-RF & 0.91 & 0.70 & 7.6 & 35.4 \\ 
\midrule\addlinespace[2.5pt]
\multicolumn{5}{l}{Stunting — MAR 50\%} \\[2.5pt] 
\midrule\addlinespace[2.5pt]
Complete data & 0.93 & 0.69 & 4.8 & 34.6 \\ 
CCA & 0.89 & 0.67 & 3.2 & 32.6 \\ 
IPW & 0.88 & 0.62 & 2.0 & 30.2 \\ 
MI-PMM & 0.90 & 0.67 & 3.4 & 35.4 \\ 
MI-RF & 0.90 & 0.72 & 8.8 & 33.0 \\ 
\midrule\addlinespace[2.5pt]
\multicolumn{5}{l}{Stunting — MAR-strong 10\%} \\[2.5pt] 
\midrule\addlinespace[2.5pt]
Complete data & 0.92 & 0.69 & 6.2 & 33.0 \\ 
CCA & 0.92 & 0.68 & 6.2 & 30.8 \\ 
IPW & 0.92 & 0.67 & 5.2 & 31.8 \\ 
MI-PMM & 0.92 & 0.67 & 4.6 & 31.6 \\ 
MI-RF & 0.92 & 0.68 & 5.8 & 31.8 \\ 
\midrule\addlinespace[2.5pt]
\multicolumn{5}{l}{Stunting — MAR-strong 30\%} \\[2.5pt] 
\midrule\addlinespace[2.5pt]
Complete data & 0.93 & 0.69 & 5.2 & 34.8 \\ 
CCA & 0.90 & 0.68 & 4.8 & 24.4 \\ 
IPW & 0.90 & 0.66 & 3.6 & 28.4 \\ 
MI-PMM & 0.90 & 0.65 & 2.8 & 23.4 \\ 
MI-RF & 0.91 & 0.70 & 6.6 & 27.0 \\ 
\midrule\addlinespace[2.5pt]
\multicolumn{5}{l}{Stunting — MAR-strong 50\%} \\[2.5pt] 
\midrule\addlinespace[2.5pt]
Complete data & 0.92 & 0.68 & 4.8 & 36.0 \\ 
CCA & 0.87 & 0.63 & 3.4 & 25.2 \\ 
IPW & 0.87 & 0.61 & 2.8 & 28.2 \\ 
MI-PMM & 0.87 & 0.59 & 1.2 & 22.6 \\ 
MI-RF & 0.89 & 0.66 & 3.6 & 27.2 \\ 
\midrule\addlinespace[2.5pt]
\multicolumn{5}{l}{Stunting — MNAR 10\%} \\[2.5pt] 
\midrule\addlinespace[2.5pt]
Complete data & 0.93 & 0.70 & 5.8 & 33.8 \\ 
CCA & 0.92 & 0.72 & 7.2 & 33.4 \\ 
IPW & 0.92 & 0.72 & 7.0 & 31.4 \\ 
MI-PMM & 0.92 & 0.72 & 7.0 & 33.2 \\ 
MI-RF & 0.93 & 0.73 & 8.0 & 33.8 \\ 
\midrule\addlinespace[2.5pt]
\multicolumn{5}{l}{Stunting — MNAR 30\%} \\[2.5pt] 
\midrule\addlinespace[2.5pt]
Complete data & 0.92 & 0.68 & 5.2 & 37.2 \\ 
CCA & 0.89 & 0.70 & 6.4 & 33.6 \\ 
IPW & 0.89 & 0.68 & 5.4 & 30.4 \\ 
MI-PMM & 0.89 & 0.69 & 5.2 & 33.6 \\ 
MI-RF & 0.90 & 0.72 & 7.6 & 31.2 \\ 
\midrule\addlinespace[2.5pt]
\multicolumn{5}{l}{Stunting — MNAR 50\%} \\[2.5pt] 
\midrule\addlinespace[2.5pt]
Complete data & 0.92 & 0.69 & 8.2 & 34.4 \\ 
CCA & 0.85 & 0.66 & 4.6 & 30.6 \\ 
IPW & 0.84 & 0.65 & 3.2 & 26.4 \\ 
MI-PMM & 0.84 & 0.66 & 2.6 & 28.8 \\ 
MI-RF & 0.87 & 0.69 & 5.6 & 27.0 \\ 
\midrule\addlinespace[2.5pt]
\multicolumn{5}{l}{Wasting — MCAR 10\%} \\[2.5pt] 
\midrule\addlinespace[2.5pt]
Complete data & 0.87 & 0.73 & 8.0 & 63.2 \\ 
CCA & 0.86 & 0.71 & 7.2 & 63.2 \\ 
IPW & 0.86 & 0.71 & 6.8 & 63.6 \\ 
MI-PMM & 0.87 & 0.72 & 6.6 & 67.0 \\ 
MI-RF & 0.86 & 0.71 & 6.6 & 64.2 \\ 
\midrule\addlinespace[2.5pt]
\multicolumn{5}{l}{Wasting — MCAR 30\%} \\[2.5pt] 
\midrule\addlinespace[2.5pt]
Complete data & 0.87 & 0.73 & 10.2 & 63.0 \\ 
CCA & 0.83 & 0.70 & 6.2 & 54.6 \\ 
IPW & 0.83 & 0.70 & 6.4 & 54.6 \\ 
MI-PMM & 0.85 & 0.72 & 7.8 & 65.0 \\ 
MI-RF & 0.84 & 0.70 & 6.0 & 61.4 \\ 
\midrule\addlinespace[2.5pt]
\multicolumn{5}{l}{Wasting — MCAR 50\%} \\[2.5pt] 
\midrule\addlinespace[2.5pt]
Complete data & 0.86 & 0.73 & 9.0 & 62.8 \\ 
CCA & 0.79 & 0.66 & 5.4 & 49.4 \\ 
IPW & 0.79 & 0.66 & 5.4 & 49.2 \\ 
MI-PMM & 0.83 & 0.70 & 7.2 & 65.4 \\ 
MI-RF & 0.80 & 0.67 & 4.8 & 60.4 \\ 
\midrule\addlinespace[2.5pt]
\multicolumn{5}{l}{Wasting — MAR 10\%} \\[2.5pt] 
\midrule\addlinespace[2.5pt]
Complete data & 0.87 & 0.72 & 7.6 & 65.0 \\ 
CCA & 0.86 & 0.71 & 6.6 & 63.2 \\ 
IPW & 0.86 & 0.71 & 6.2 & 60.8 \\ 
MI-PMM & 0.87 & 0.72 & 7.2 & 64.2 \\ 
MI-RF & 0.87 & 0.72 & 7.4 & 59.6 \\ 
\midrule\addlinespace[2.5pt]
\multicolumn{5}{l}{Wasting — MAR 30\%} \\[2.5pt] 
\midrule\addlinespace[2.5pt]
Complete data & 0.87 & 0.73 & 8.4 & 64.8 \\ 
CCA & 0.84 & 0.68 & 4.0 & 55.8 \\ 
IPW & 0.83 & 0.69 & 4.4 & 52.4 \\ 
MI-PMM & 0.86 & 0.72 & 7.6 & 57.2 \\ 
MI-RF & 0.85 & 0.71 & 6.2 & 49.6 \\ 
\midrule\addlinespace[2.5pt]
\multicolumn{5}{l}{Wasting — MAR 50\%} \\[2.5pt] 
\midrule\addlinespace[2.5pt]
Complete data & 0.87 & 0.72 & 9.0 & 63.2 \\ 
CCA & 0.80 & 0.65 & 4.4 & 52.0 \\ 
IPW & 0.78 & 0.65 & 4.4 & 48.4 \\ 
MI-PMM & 0.84 & 0.69 & 7.6 & 56.4 \\ 
MI-RF & 0.82 & 0.68 & 4.8 & 40.0 \\ 
\midrule\addlinespace[2.5pt]
\multicolumn{5}{l}{Wasting — MAR-strong 10\%} \\[2.5pt] 
\midrule\addlinespace[2.5pt]
Complete data & 0.87 & 0.72 & 7.0 & 63.0 \\ 
CCA & 0.86 & 0.68 & 4.4 & 58.8 \\ 
IPW & 0.86 & 0.71 & 6.4 & 60.0 \\ 
MI-PMM & 0.87 & 0.70 & 5.4 & 61.6 \\ 
MI-RF & 0.86 & 0.69 & 5.4 & 59.8 \\ 
\midrule\addlinespace[2.5pt]
\multicolumn{5}{l}{Wasting — MAR-strong 30\%} \\[2.5pt] 
\midrule\addlinespace[2.5pt]
Complete data & 0.87 & 0.73 & 8.8 & 64.8 \\ 
CCA & 0.82 & 0.60 & 1.0 & 46.0 \\ 
IPW & 0.82 & 0.66 & 3.0 & 50.2 \\ 
MI-PMM & 0.85 & 0.67 & 3.0 & 53.6 \\ 
MI-RF & 0.82 & 0.62 & 2.4 & 49.6 \\ 
\midrule\addlinespace[2.5pt]
\multicolumn{5}{l}{Wasting — MAR-strong 50\%} \\[2.5pt] 
\midrule\addlinespace[2.5pt]
Complete data & 0.87 & 0.72 & 6.2 & 61.6 \\ 
CCA & 0.76 & 0.55 & 0.4 & 37.0 \\ 
IPW & 0.76 & 0.60 & 1.2 & 41.4 \\ 
MI-PMM & 0.83 & 0.64 & 2.6 & 46.4 \\ 
MI-RF & 0.77 & 0.56 & 0.4 & 40.2 \\ 
\midrule\addlinespace[2.5pt]
\multicolumn{5}{l}{Wasting — MNAR 10\%} \\[2.5pt] 
\midrule\addlinespace[2.5pt]
Complete data & 0.87 & 0.72 & 8.0 & 64.4 \\ 
CCA & 0.84 & 0.72 & 9.2 & 23.2 \\ 
IPW & 0.84 & 0.72 & 10.2 & 19.0 \\ 
MI-PMM & 0.85 & 0.74 & 11.6 & 21.4 \\ 
MI-RF & 0.84 & 0.72 & 8.6 & 19.6 \\ 
\midrule\addlinespace[2.5pt]
\multicolumn{5}{l}{Wasting — MNAR 30\%} \\[2.5pt] 
\midrule\addlinespace[2.5pt]
Complete data & 0.87 & 0.73 & 7.8 & 61.2 \\ 
CCA & 0.77 & 0.69 & 8.8 & 11.4 \\ 
IPW & 0.76 & 0.70 & 9.4 & 8.4 \\ 
MI-PMM & 0.81 & 0.73 & 10.8 & 10.4 \\ 
MI-RF & 0.78 & 0.70 & 7.6 & 6.4 \\ 
\midrule\addlinespace[2.5pt]
\multicolumn{5}{l}{Wasting — MNAR 50\%} \\[2.5pt] 
\midrule\addlinespace[2.5pt]
Complete data & 0.87 & 0.73 & 9.4 & 62.2 \\ 
CCA & 0.67 & 0.63 & 3.0 & 16.0 \\ 
IPW & 0.66 & 0.65 & 2.8 & 12.2 \\ 
MI-PMM & 0.76 & 0.71 & 6.0 & 18.8 \\ 
MI-RF & 0.70 & 0.66 & 3.6 & 8.6 \\ 
\bottomrule
\end{longtable}
\endgroup

\clearpage

% --- tableS5_ntree_substudy
\begingroup
\fontsize{8.0pt}{10.0pt}\selectfont
\begin{longtable}{l|rrrrrr}
\caption*{
{\fontsize{20}{25}\selectfont  Table S5. MI-RF with 10 versus 50 trees (MAR, 30\% missing)\fontsize{8}{10}\selectfont }
} \\ 
\toprule
 & R & Bias, pp (MCSE) & Empirical SE & Model SE & SE ratio & Coverage \% (MCSE) \\ 
\midrule\endhead\addlinespace[2.5pt]
\multicolumn{7}{l}{Stunting} \\[2.5pt] 
\midrule\addlinespace[2.5pt]
MI-RF & 500 & -0.09 (0.05) & 1.04 & 0.97 & 0.93 & 93.4 (1.1) \\ 
MI-RF (ntree=50) & 200 & -0.10 (0.07) & 0.97 & 0.97 & 0.99 & 96.5 (1.3) \\ 
\midrule\addlinespace[2.5pt]
\multicolumn{7}{l}{Wasting} \\[2.5pt] 
\midrule\addlinespace[2.5pt]
MI-RF & 500 & -0.15 (0.02) & 0.48 & 0.48 & 1.00 & 94.0 (1.1) \\ 
MI-RF (ntree=50) & 200 & -0.18 (0.03) & 0.44 & 0.48 & 1.09 & 97.0 (1.2) \\ 
\bottomrule
\end{longtable}
\endgroup

\clearpage

\section*{Figures}

\begin{figure}[p]
\centering
\includegraphics[width=\linewidth,height=0.86\textheight,keepaspectratio]{D:/project/handling missing data/outputs/figures/figure1_equity_ranking.pdf}
\caption{\textbf{Subgroup prevalence and equity ranking under each missing-data method.} Prevalence of stunting and wasting by region and wealth quintile under each missing-data method. Grey bars: 95\% confidence interval for complete-case analysis (CCA). Points: point estimates under each method. Groups are ordered by CCA prevalence; the dashed line marks the five worst-affected regions under CCA. Daggers mark subgroups whose top-5 worst-affected status differs between methods. Panel subtitles give the minimum Spearman correlation and top-k overlap across all method pairs, and the sampling-noise benchmark: the median agreement between the CCA point ranking and rankings re-drawn from the CCA sampling distribution (10,000 draws).}
\end{figure}
\clearpage

\begin{figure}[p]
\centering
\includegraphics[width=\linewidth,height=0.86\textheight,keepaspectratio]{D:/project/handling missing data/outputs/figures/figure2_simulation_coverage.pdf}
\caption{\textbf{Confidence-interval coverage by missingness mechanism and rate.} Coverage of nominal 95\% confidence intervals of survey-weighted national prevalence of stunting (top row) and wasting (bottom row) by missingness mechanism (columns) and missingness rate, from 500 simulation replicates per scenario. Error bars: plus or minus 1.96 Monte Carlo standard errors. Complete data: the estimate from each bootstrap sample before deletion (benchmark). The shaded band is the range expected under nominal coverage (95\% plus or minus 1.96 Monte Carlo standard errors).}
\end{figure}
\clearpage

\begin{figure}[p]
\centering
\includegraphics[width=\linewidth,height=0.86\textheight,keepaspectratio]{D:/project/handling missing data/outputs/figures/figureS1_equity_ranking_sensitivity.pdf}
\caption{\textbf{Subgroup prevalence and equity ranking, sensitivity analysis.} Prevalence of stunting and wasting by region and wealth quintile under each missing-data method, excluding children with WHO-flagged z-scores from the target population. Grey bars: 95\% confidence interval for complete-case analysis (CCA). Points: point estimates under each method. Groups are ordered by CCA prevalence; the dashed line marks the five worst-affected regions under CCA. Daggers mark subgroups whose top-5 worst-affected status differs between methods. Panel subtitles give the minimum Spearman correlation and top-k overlap across all method pairs, and the sampling-noise benchmark: the median agreement between the CCA point ranking and rankings re-drawn from the CCA sampling distribution (10,000 draws).}
\end{figure}
\clearpage

\begin{figure}[p]
\centering
\includegraphics[width=\linewidth,height=0.86\textheight,keepaspectratio]{D:/project/handling missing data/outputs/figures/figureS2_simulation_bias.pdf}
\caption{\textbf{Bias of national prevalence by missingness mechanism and rate.} Bias, in percentage points, of survey-weighted national prevalence of stunting (top row) and wasting (bottom row) by missingness mechanism (columns) and missingness rate, from 500 simulation replicates per scenario. Error bars: plus or minus 1.96 Monte Carlo standard errors. Complete data: the estimate from each bootstrap sample before deletion (benchmark). Vertical scales differ between outcomes.}
\end{figure}
\clearpage

\end{document}
